% concatenated from 05_25_2025_submit_Sharpness_of_L2_Fourier_restriction_in_higher_dimensions.tex
\documentclass{amsart}
%Use the enumitem package instead of enumerate

\usepackage[normalem]{ulem}

\usepackage[shortlabels]{enumitem}
\renewcommand{\descriptionlabel}[1]{%
\hspace\labelsep \upshape\bfseries #1%
}
%\usepackage{enumitem}

%then it will support the same suntax as the enumerate package.
%The enumerate package does not provide any extra configurations other than the label.
\numberwithin{equation}{section}

\usepackage{amssymb,amsmath,amsthm,verbatim,IEEEtrantools,hyperref}

\theoremstyle{definition}
\newtheorem{theorem}{Theorem}[section]
\newtheorem{thm}[theorem]{Theorem}
\newtheorem*{lemma*}{Lemma}
\newtheorem{lemma}[theorem]{Lemma}
\newtheorem{corollary}[theorem]{Corollary}
\newtheorem{prop}[theorem]{Proposition}
\newtheorem{example}[theorem]{Example}
\newtheorem{remark}[theorem]{Remark}
\newtheorem{defn}[theorem]{Definition}


\theoremstyle{definition}
\newtheorem*{notation}{Notation}
\newtheorem*{Outline-of-the-Rest-of-the-Paper}{Outline of the Rest of the Paper}

\def\forinfmany{\text{ for infinitely many }}
\def\P{P}

\def\RR{\mathbb{R}}
\def\CC{\mathbb{C}}
\def\FF{\mathbb{F}}
\def\QQ{\mathbb{Q}}
\def\ZZ{\mathbb{Z}}
\def\NN{\mathbb{N}}
%\def\NN{\mathbb{Z}_{> 0}}
\def\Nzero{\mathbb{Z}_{\geq 0}}
\def\EE{\mathbb{E}}
\def\PP{\mathbb{P}}
\def\calC{\mathcal{C}}
\def\calD{\mathcal{D}}
\def\calE{\mathcal{E}}
\def\calF{\mathcal{F}}
\def\calG{\mathcal{G}}
\def\calH{\mathcal{H}}
\def\supp{\mathrm{supp}}
\def\diam{\mathrm{diam}}
\def\sp{\mathrm{span}}
\def\ker{\mathrm{ker}}
%\def\sp{\mathrm{span}} %messes up align enviroment
\newcommand{\rbr}[1]{\left( {#1} \right)}
\newcommand{\sbr}[1]{\left[ {#1} \right]}
\newcommand{\cbr}[1]{ \left\{ {#1} \right\} }
\newcommand{\abr}[1]{\left\langle {#1} \right\rangle}
\newcommand{\abs}[1]{\left| {#1} \right|}
\newcommand{\norm}[1]{\left\|#1\right\|}
\def\one{\mathbf{1}}
\DeclareMathOperator*{\esssup}{ess\,sup}
\newcommand*\wc{{}\cdot{}}
%\newcommand*\wc{ \, \cdot \,}
%wc for wildcard
\renewcommand{\Re}{\operatorname{Re}}
\renewcommand{\Im}{\operatorname{Im}}
\newcommand{\OK}{\mathcal{O}_K}
\newcommand{\isom}{\cong}

\newcommand{\prin}{\mathrm{prin}}
\def\con{\tilde{C}}
\DeclareMathOperator{\tr}{tr}
\DeclareMathOperator{\dist}{dist}

%\usepackage{lineno}
%\linenumbers 

%\usepackage[notref,notcite]{showkeys}


\usepackage{xcolor}
\def\donggeun#1{\noindent
\textcolor{blue}
{\textsc{(Donggeun:}
\textsf{#1})}}

\def\kyle#1{\noindent
\textcolor{purple}
{\textsc{(Kyle:}
\textsf{#1})}}


\def\robert#1{\noindent
\textcolor{red}
{\textsc{(Robert:}
\textsf{#1})}}



\begin{document}

\author{Robert Fraser \and Kyle Hambrook \and Donggeun Ryou}
\title[Sharpness of the Mockenhaupt-Mitsis-Bak-Seeger theorem]{Sharpness of the Mockenhaupt-Mitsis-Bak-Seeger  Fourier restriction theorem in all dimensions}

\thanks{This material is based upon work supported by the National Science Foundation under
Award No. 2400329}

\begin{abstract}
We prove the optimality of the exponent in the Mockenhaupt-Mitsis-Bak-Seeger Fourier restriction theorem in all dimensions $d$ and the full parameter range $0 < a,b < d$. Our construction is  deterministic and also yields Salem sets. 
\end{abstract}

\maketitle

\section{Introduction}
\subsection{Main Result}
This article concerns the optimality of 
the constant $p_{*}(a,b,d)$ in 
the following general Fourier restriction theorem. 








\begin{thm}[Mockenhaupt-Mitsis-Bak-Seeger]\label{frac ST}
Let $0 < a, b < d$, with $d$ an integer. 
Define $p_{*}(a,b,d) = (4d-4a+2b)/b.$
Suppose $\mu$ is a Borel  measure on $\RR^d$ with the following properties: 
\begin{enumerate}[(A)]
    \item $\mu(B(x,r)) 
    %\leq C 
\lesssim 
    r^{a}$ for all $x \in \RR^d$ and all $r > 0$.  
    \item $|\widehat{\mu}(\xi)| \lesssim 
    %\leq C 
    (1+|\xi|)^{-b/2}$ for all $\xi \in \RR^d$.  
\end{enumerate}
Then, for each $p \geq p_{*}(a,b,d)$, the following holds: 
\begin{enumerate}[(A)]
\setcounter{enumi}{17}
\item $\|\widehat{f\mu}\|_{L^p} 
%\leq C_p 
\lesssim_p 
\|f\|_{L^2(\mu)}$ for all $f \in L^2(\mu)$.
\end{enumerate}
\end{thm}

\noindent The notation is explained in Subsection \ref{notation}. The history of this theorem and of the question of the optimality of the constant $p_{*}(a,b,d)$ is deferred to Section \ref{motivation}. 

The main result of this article is that the constant $p_{*}(a,b,d)$ in Theorem \ref{frac ST} is best possible for all values of $0 < a,b < d$. 

We prove our main result by constructing a measure   
on a set of vectors that are well-approximable by elements of lattices arising from algebraic number theory. 

Let $K$ be a number field (i.e., a finite extension field of $\QQ$) of degree $d$. 
Let $\OK$ be the ring of integers of $K$.
For each ideal $I \subseteq \OK$, let $N(I)$ denote the norm of $I$ and $I^{-1}$ denote the fractional ideal which is the inverse of $I$. 
Let $B=\cbr{\omega_1,\ldots,\omega_{d}}$ be an integral basis for $K$ (i.e., a basis for $K$ as a $\QQ$-vector space which is also a basis for $\OK$ as a $\ZZ$-module). We identify $\mathbb{Q}^d$ with $K$ by identifying 
$q=(q_1,\ldots,q_n) \in \mathbb{Q}^d$ with 
$q=\sum_{i=1}^{d} q_i \omega_i \in K$. 
Since $B$ is an integral basis, 
this also identifies $\mathbb{Z}^d$ with $\OK$. 
Moreover, if $I \subseteq \OK$ is an integral ideal, $I^{-1}$ is identified with a lattice containing $\mathbb{Z}^d$. 
This identification also extends to lengths and distances:  
If $q=\sum_{i=1}^d q_i \omega_i \in K$ 
and $x = (x_1,\ldots,x_d) \in \RR^d$, 
then 
%$|q| = \max_{1 \leq i \leq n} |q_i|$. 
$|q| = \rbr{\sum_{i=1}^d |q_i|^2}^{1/2}$  
and $|q-x| = \rbr{\sum_{i=1}^d |q_i-x_i|^2}^{1/2}$. 


Let $\tau > 1$. Define 
$E(K,B,\tau)$ to be the set of all $x \in \RR^d$ such that 
$$
\dist(x, I^{-1}) \leq |N(I)|^{-(\tau+1)/d} 
$$
for infinitely many ideals $I$ of $\OK$.
Here 
$
\dist(x,I^{-1}) = \min\cbr{|x-t|: t \in I^{-1}}. 
$


\begin{thm}[Main Theorem]\label{mainthm3}
Let $1 \leq q < \infty$. 
Let $-d < \rho < d$ 
and define  
$
\rho^+ = \max\cbr{\rho,0}$ and $\rho^- = \max\cbr{-\rho,0}
$. 
There exists a Borel probability measure 
$\mu$ with the following properties: 
\begin{enumerate}[(i)]
    \item $\supp(\mu) \subseteq E(K, B, \tau)$; 
\item $\mu$ satisfies (A) for each exponent $a < \frac{2d  - \rho^-}{1 + \tau}$;
\item 
$\mu$ satisfies (B) for each  exponent $b < \frac{2(d-\rho^+)}{1 + \tau}$; 
\item There exists a sequence of functions $f_k \in L^q(\mu)$ such that
\[\lim_{k \to \infty} \frac{\|\widehat{f_k \mu}\|_{L^p}}{\|f_k\|_{L^q(\mu)}} = \infty\]
whenever $p < p(\tau, \rho, q, d) := \frac{q(d \tau - \rho)}{(q-1) (d-\rho^+)}$.
\end{enumerate}
\end{thm}


Subsection \ref{ideas} describes the main ideas and novel features of the proof of this theorem.  
%Theorem \ref{mainthm3}. 
The formal proof of 
the theorem 
%Theorem \ref{mainthm3} 
constitutes Sections \ref{Algebraic Number Theory: Basics section} to \ref{sec_Failure_Estimate}.

Theorem \ref{mainthm3} and the following remark 
show that, for any integer $d$ and any real numbers $a$ and $b$ satisfying $0 < a, b < d$, the exponent $p_*(a,b,d)$ in  the Mockenhaupt-Mitsis-Bak-Seeger restriction theorem cannot be improved.



\begin{remark}
We may assume $b \leq 2a$. Indeed, Mitsis \cite[Proposition 3.1]{mitsis-res} 
showed that if (A) and (R) hold, then 
$p \geq 2d/a$.
By combining this with Theorem \ref{frac ST}, we see that (A) and (B) imply $b \leq 2a$. 
Let $p_0 < p_{*}(a,b,d)$. Choose $a_0$ and $b_0$ such that $0 < a < a_0 < d$, $0 < b < b_0 < d$, $b_0 \leq 2a_0$, and $p_0 < p_{*}(a_0,b_0,d) < p_{*}(a,b,d)$. 
If $b_0 \leq a_0$, choose $\tau = \frac{2d}{a_0} - 1$ and $\rho = d\left(1-\frac{b_0}{a_0}\right)$. 
If $a_0 \leq b_0 \leq 2a_0$, choose $\tau = \frac{2d}{b_0} - 1$ and $\rho = 2d\left(\frac{a_0}{b_0}-1\right)$. 
In both cases, $a_0 = \frac{2d  - \rho^-}{1 + \tau}$,  $b_0 = \frac{2(d  - \rho^+)}{1 + \tau}$, and 
$p_{*}(a_0,b_0,d) = p(\tau,\rho,2,d)$. 
Then Theorem \ref{mainthm3} gives a measure $\mu$ satisfying (A) and (B) but failing (R) for $p=p_0$ (and every other $p < p_{*}(a_0,b_0,d)$). 
\end{remark}


\subsection{Deterministic Salem Sets and a Variant of \texorpdfstring{$E(K,B,\tau)$}{E(k,B,tau}.} 

A set in $\RR^d$ is called a Salem set if it supports a Borel measure $\mu$ that satisfies (B) for every value of $b$ less than the Hausdorff dimension of the set. 
There are many random constructions of Salem sets in $\RR^d$ (see \cite{salem}, \cite{kahane-book}, \cite{kahane-1966-fourier}, \cite{kahane-1966-brownian}, \cite{bluhm-1}, \cite{chen-seeger}, \cite{ekstrom}, \cite{LP}, \cite{shmerkin-suomala}). 
The first deterministic Salem sets in $\RR$ of arbitrary dimension were found by 
Kaufman \cite{Kaufman}, who proved 
$$
E(\tau) = \cbr{x \in \RR : |x-r/q| \leq |q|^{-(1+\tau)} \forinfmany (q,r) \in \ZZ \times \ZZ }
$$
is a Salem set of dimension $2/(1+\tau)$ when $\tau > 1$. 
 


Fraser and Hambrook \cite{hambrook-fraser-Rn}  
extended Kaufman's result to $\RR^d$ 
by showing the set 
$$
E_{\prin}(K,B,\tau)=
\cbr{x \in \RR^d : |x-r/q| \leq |q|^{-(\tau+1)} \text{ for infinitely many } (q,r) \in \OK \times \OK}
$$
is a Salem set of dimension $2d/(1+\tau)$ when $\tau > 1$. 


The set $E_{\prin}(K,B,\tau)$ is a variant of $E(K,B,\tau)$ defined in terms of \textit{elements} rather than \textit{ideals}. 
In the present article, we primarily considered $E(K, B, \tau)$ rather than $E_{\prin}(K, B, \tau)$ because working with \textit{ideals} rather than \textit{elements} simplifies many parts of the argument. However, the two sets have many of the same properties, as the next two results illustrate. 

By Frostman's lemma, the exponent $b$ in (B) can never exceed the Hausdorff dimension of the support of $\mu$ (see, e.g.,  \cite[Section 3.6]{mattila-book-2}). 
Thus, Theorem \ref{mainthm3} combined with Proposition \ref{hausdorff upper bound} (which gives an upper bound on the Hausdorff dimension of $E(K,B,\tau)$)
implies 


\begin{corollary}\label{Salem_cor}
$E(K,B,\tau)$ is a Salem set of dimension $2d/(1+\tau)$ for every $\tau > 1$.
\end{corollary}













By modifying the proof of Theorem \ref{mainthm3} as described in Section \ref{sec_prin}, we can prove 

%
\begin{thm}\label{Eprinthm}
The statement of Theorem \ref{mainthm3} also holds for $E_{\prin}(K, B, \tau)$.
\end{thm}
%Theorem \ref{mainthm3} can also be shown to hold for $E_{\prin}(K, B, \tau)$. 
%We explain the necessary modifications to the proof of Theorem \ref{mainthm3} in Section \ref{sec_prin}. 








\subsection{Notation}\label{notation}
For a measure $\nu$ on $\RR^d$, the Fourier transform of $\nu$ is   $$\widehat{\nu}(\xi)=\int_{\RR^d} e^{-2\pi i x \cdot \xi } d\nu(x) \quad \text{for all $\xi \in \RR^d$}.$$
We make the following exception to 
this notation: If $F:\RR^d \to \CC$ is $\ZZ^d$-periodic, 
$$
\widehat{F}(s)  
=\int_{[0,1)^d} e^{-2\pi i x \cdot s } F(x) dx \quad \text{for all $s \in \ZZ^d$}.$$


For a measure $\nu$ on $\RR^d$, $L^q(\nu)$ is the space of all measurable function $f:\RR^d \to \RR$ such that $\int |f|^q d\nu < \infty$, with functions that are equal $\nu$-a.e.\ identified.
If $\nu$ is Lebesgue measure, we write $L^q$ instead of $L^q(\nu)$.

We use $B(x,r)$  to denote the closed ball with center $x$ and radius $r$ in the Euclidean norm on $\RR^d$. 


For any set $A$, the indicator function of $A$ is the function $\one_A$  defined by $\one_A(x) = 1$ if $x \in A$ and $\one_A(x) = 0$ otherwise. 

A statement of the form ``$F(u) \lesssim G(u)$ for all $u$'' means ``there exists a constant $C>0$ such that $F(u) \leq C G(u)$ for all $u$.'' If the implied constant $C$ depends on a relevant parameter, say $p$, we write $\lesssim_{p}$ instead of $\lesssim$. 
Note also 
$F(u)  \gtrsim G(u)$ means $G(u) \lesssim F(u)$,   
and 
$F(u) \sim G(u)$ means $F(u) \lesssim G(u)$ and $F(u) \gtrsim  G(u)$.


\subsection{Organization of The Paper}\label{organization}

Section \ref{motivation} contains motivation for Theorem \ref{mainthm3} and discussion about its proof. 

The proof of Theorem \ref{mainthm3} constitutes sections \ref{Algebraic Number Theory: Basics section} to \ref{sec_Failure_Estimate}.
In Section \ref{Algebraic Number Theory: Basics section}, we introduce the necessary algebraic number theory we need. 
In Section \ref{sec_Single-
Scale}, we give the definitions and prove Fourier estimates for some functions needed for the construction of the measure $\mu$. 
In Section \ref{Intermediate Measures}, we introduce some intermediate measures and prove Fourier-analytic estimates  that will be used in later sections. 
In Section \ref{sec_measure}, we define the measure $\mu$ and show that it satisfies the support and Fourier decay conditions in Theorem \ref{mainthm3}. In particular,  
Proposition \ref{support of mu lemma} and Proposition \ref{Prop_mu_decay}, respectively,  establish (i) and (iii) in Theorem \ref{mainthm3}. 
In Section \ref{regularity_section}, we prove the regularity property in Theorem \ref{mainthm3}. In particular, Proposition \ref{Prop_mu_reg} establishes (ii) in Theorem \ref{mainthm3}.  
In Section \ref{sec_Failure_Estimate}, we prove $\mu$ fails the restriction property (R) for the appropriate values of $p$. In particular,  Proposition \ref{res fail prop} establishes (iv) in Theorem \ref{mainthm3}. 


Section \ref{convstab section} contains the Convolution Stability Lemma (Lemma \ref{convstab lemma}),  which is used in Section \ref{sec_measure} to prove necessary Fourier decay estimates. We have separated  this lemma from the main proof because it applies more  generally. 
In Section \ref{Hausdorff Dimension Upper Bound section}, we prove an upper bound on the Hausdorff dimension of $E(K,B,\tau)$. 
Section \ref{sec_prin} contains the modifications to the proof of Theorem \ref{mainthm3} necessary to prove Theorem \ref{Eprinthm}.


\section{Motivation and Discussion}\label{motivation}

\subsection{Restriction Theorems}\label{restriction theorems}

The problem of restriction of the Fourier transform to small sets is an interesting one with a long history (which we describe only briefly). Because the Fourier transform is an isometry on $L^2$, there is no sensible way to restrict the Fourier transform of an arbitrary $L^2$ function to a set of Lebesgue measure zero. 
On the other hand, because the Fourier transform of a function in $L^1$ is continuous, 
the Fourier transform of an $L^1$ function can be sensibly restricted to any measurable subset of $\mathbb{R}^d$. The area of \textit{Fourier restriction} concerns ways of restricting the Fourier transform of  
functions in spaces between  $L^1$ and $L^2$
to certain sets of Lebesgue measure zero. 


We say that the measure $\mu$ on $\RR^d$  satisfies a \textit{Fourier extension} estimate with exponents $p$ and $q$ if
\begin{equation}\label{restriction}
\norm{\widehat{f \mu}}_{L^p} \lesssim_{p,q} 
%\leq C
\norm{f}_{L^q(\mu)}
\text{  for all $f \in L^q(\mu)$.}
\end{equation}
%
It can be seen using a duality argument (see, e.g., \cite{mattila-book-2}) that the estimate \eqref{restriction} implies that the Fourier transform of functions in $L^{p'}$ can be meaningfully restricted to the support of the measure $\mu$. More precisely, \eqref{restriction} is equivalent to the statement that the Fourier transform map $f \mapsto \widehat{f}$ is a bounded linear map from $L^{p'}$ to $L^{q'}(\mu)$. 
Here $p'$ and $q'$ are the conjugate exponents to $p$ and $q$, respectively. 

The foundational result in Fourier restriction theory is the Stein-Tomas theorem \cite{tomas-1975}, \cite{tomas-symp},  which says that \eqref{restriction} holds for $q=2$ and $p \geq 2(d+1)/(d-1)$ when $\mu$ is the surface measure on a sphere in $\RR^d$. Notable generalizations of the Stein-Tomas theorem to the surface measure on quadratic hypersurfaces and arbitrary smooth submanifolds in $\RR^d$ were obtained, respectively, by Strichartz \cite{Strichartz77} and Greenleaf \cite{Greenleaf}. 

All of the restriction theorems mentioned in the precious paragraph are special cases of 
Theorem \ref{frac ST}, the \textit{Mockenhaupt-Mitsis-Bak-Seeger restriction theorem}. 
Mockenhaupt \cite{mock} and Mitsis \cite{mitsis-res} independently proved Theorem \ref{frac ST} with $p > p_{*}(a,b,d)$;  
the endpoint case $p = p_{*}(a,b,d)$ was established by Bak and Seeger \cite{bak-seeger}. 



Condition (A) in Theorem \ref{frac ST} is a well-known condition in geometric measure theory sometimes called the Frostman condition or ball condition. 
Mitsis \cite{mitsis-res} proved that condition (B) in Theorem \ref{frac ST} implies condition (A) with $a=b/2$. 
This makes Theorem \ref{frac ST} quite remarkable ---  it allows one use the  pointwise estimate (B) on $\widehat{\mu}$ to conclude an $L^p$ decay estimate (R) on $\widehat{f\mu}$ for all functions $f \in L^2(\mu)$. 



 

\subsection{Sharpness of Theorem \ref{frac ST}.}


Mockenhaupt\cite{mock} and Mitsis \cite{mitsis-res} 
both noted the question of whether the constant $p_*(a,b,d)$ in Theorem \ref{frac ST} is best possible was an open problem. 









When $d \geq 2$ and $a=b=d-1$, the constant $p_{*}(a,b,d)$ in Theorem \ref{frac ST} is known to be optimal because of an example due to Knapp (according to  \cite{tomas-symp}) for the surface measure on the sphere. 
The original purpose of Knapp's example was to show that the constant $2(d+1)/(d-1)$ in the Stein-Tomas theorem (mentioned in Subsection \ref{restriction theorems}) is best possible. 
The Knapp example was later extended 
 to smooth hypersurfaces with nonvanishing Gaussian curvature  by Strichartz \cite{Strichartz77}. See also Example 1.8 of Demeter \cite{Demeter} for a further generalization of Knapp's example  that shows \eqref{restriction} cannot hold if $\mu$ is a measure on a smooth hypersurface and $p < \frac{q(d + 1)}{(q - 1) (d - 1)}.$ 

 Hambrook and {\L}aba \cite{HL2013}  proved the optimality of $p_*(a, b, d)$ in Theorem \ref{frac ST} when $d=1$ and $0 < b \leq a \leq d$. 
 (In fact, Hambrook and {\L}aba \cite{HL2013} only claimed the case $a=b$ for $a,b$ of a special form.   
Hambrook \cite{hambrook-thesis} 
showed how optimality for all $b \leq a$ could be deduced from the result of \cite{HL2013}.  
Chen \cite{chen} modified the construction of 
\cite{HL2013} to obtain a technically stronger result that also implied optimality for $b \leq a$.) 
By combining 
the ideas of Knapp's example and \cite{HL2013}, 
Hambrook and {\L}aba \cite{HL2016} proved the optimality of $p_*(a,b,d)$ for $d \geq 2$ and $d-1 \leq b \leq a \leq d$.  
Note that  none of the sharpness results mentioned above covered the case $b>a$. Fraser, Hambrook, and Ryou 
\cite{fraser-hambrook-ryou} 
built upon the idea  of  \cite{HL2013} to
prove the optimality of $p_*(a,b,d)$ when $d=1$ in the full range $0 < a,b < d$ (notably, the construction is  deterministic). 
Li and Liu \cite{bochen-dimension} independently obtained the same optimality result for $d=1$ by a similar construction. 
The present article 
fully resolves 
the problem of the optimality of $p_*(a,b,d)$ for \textit{all} $d \geq 1$ and the full range $0 < a,b < d$. 



\begin{remark}
We have considered the question of the sharpness of $p_*(a,b,d)$ in Theorem \ref{frac ST} by fixing $a,b,d$ and asking  whether the statement holds with $p_*(a,b,d)$ replaced by a smaller number. An alternative way to consider the question  is to look for measures which satisfy (A), (B), and (R) for $p < p_{*}(a,b,d)$.  For some results in this vein, see Shmerkin and Suomala \cite{shmerkin-suomala}, Chen and Seeger \cite{chen-seeger}, {\L}aba and Wang \cite{laba-wang}, and Ryou \cite{Ryou-2023}. 
\end{remark}


\subsection{Key Ideas in the Proof of Theorem \ref{mainthm3}}\label{ideas}

The key to failure of restriction estimates in Knapp's example is \textit{constructive interference}. 
The smoothness of hypersurfaces guarantees that a $\delta$-neighborhood of a point on the hypersurface will be contained in a $\delta \times \cdots \times \delta \times \delta^2$-rectangle oriented along the tangent plane to the hypersurface. 
The uncertainty principle (or a direct calculation) implies that 
the Fourier transform of the indicator function of this $\delta$-ball should therefore be $\gtrsim \delta^{d-1}$ in a dual rectangle with side lengths $\delta^{-1} \times \cdots \times \delta^{-1} \times \delta^{-2}$. 

When studying restriction problems for \textit{fractal sets}, the easiest way to achieve constructive interference is to choose the indicator function of a neighborhood of an arithmetic progression. For simplicity, we will consider the case $d = 1$. In fact, the indicator function of a $\delta$-neighborhood of an arithmetic progression with common difference $z$ will have large Fourier transform near frequencies of the form $\xi = n z, n \in \mathbb{Z}, |n| \leq \delta^{-1}$. This observation is used by Hambrook and {\L}aba \cite{HL2013} and a later work by Chen \cite{chen} to provide a random Cantor measure for which the exponent in (R) cannot be improved beyond $p_*(a,b,d)$.

In order to construct deterministc examples of fractal sets allowing for constructive interference, 
it is useful to have an example of a set supporting measures satisfying (A) and (B) for 
prescribed 
exponents $a$ and $b$ that also contains long arithmetic progressions. 
A natural choice for this set is the set $E(\tau)$ provided by Kaufman \cite{Kaufman}. This set is a good choice for two reasons: 
first, it is easy to locate a compact subset $F$ of $E(\tau)$ for which fattenings of $F$ contain long arithmetic progressions; 
second, the limsup structure of the set $E(\tau)$ allows for the insertion of arithmetic progressions of different length if necessary (it turns out to be necessary in the case where $a \neq b$). 
These observations are central to the previous work of Fraser, Hambrook and Ryou \cite{fraser-hambrook-ryou} 
locating measures on $E(\tau)$ for which the exponent appearing in (R) cannot be improved beyond $p_*(a,b,d)$. 

In the same spirit, it is natural to seek measures on the set $E(K,B,\tau)$ as possible counterexamples to (R) for $p < p_*(a,b,d)$. 
First, the set $E(K,B,\tau)$ 
supports measures satisfying (A) and (B) for prescribed  $a$ and $b$; 
second, it is easy to locate a compact subset of $E(K,B,\tau)$ is naturally concentrated around a small neighborhood of a small number of lattices; third, because $E(K,B,\tau)$ is a limsup set, it is easy to insert a generalized arithmetic progression of a desired length in order to obtain a counterexample if $a \neq b$.

Kaufman \cite{Kaufman} constructed a measure on $E(\tau)$ of maximal Fourier decay. In 
\cite{fraser-hambrook-ryou}, Fraser, Hambrook, and Ryou 
modified Kaufman's construction (and Papadimitropoulos's construction \cite{pa-thesis}) so that the measure satisfied (A) and (B) for prescribed values of $a$ and $b$ and failed $(R)$ for $p<p_*(a,b,d)$. As alluded to above, inserting (and modifying the weight on) a generalized arithmetic progression is the key to establishing the failure of (R). 

Hambrook and Fraser \cite{hambrook-fraser-Rn} constructed a measure on 
$E_{\prin}(K,B,\tau)$ 
with  maximal Fourier decay. 
The proof of Theorem \ref{mainthm3} 
modifies 
the Hambrook and Fraser construction 
using the ideas of 
\cite{fraser-hambrook-ryou}. 
The method used to establish the necessary Fourier decay estimate in Theorem \ref{mainthm3}  is essentially the same as the method in  \cite{hambrook-fraser-Rn}. 
However, the set $E(K,B,\tau)$ in Theorem \ref{mainthm3} is defined using ideals while the set $E_{\prin}(K,B,\tau)$ in \cite{hambrook-fraser-Rn} is defined using elements. 
This change streamlines the argument in the present paper in some ways (less explicit reliance on the basis and a simpler proof of the divisor bound). 
However, it required some additional technology (the notion of the different ideal of a number field and a more abstract exponential sum calculation). 
The regularity property in Theorem \ref{mainthm3} was not considered in \cite{hambrook-fraser-Rn}. 
The argument in the present paper generalizes the argument of Papadimitropoulos \cite{pa-thesis} to the number field setting. 
The key ingredient is the separation lemma (Lemma \ref{separation lemma}) and its application in Section \ref{regularity_section}. 
As \cite{fraser-hambrook-ryou} (and in contrast to \cite{pa-thesis}) we use a Fourier-analytic argument in proof of regularity (see Step 2 in the proof of Lemma \ref{k-ball regularity}). 
The failure of the estimate (R) is proved by an argument similar to that used in \cite{fraser-hambrook-ryou}. 
However, the argument requires more advanced ideas due to the more general setting. The role of the arithmetic progression at each scale is played by the inverse of a prime ideal. 
As in \cite{fraser-hambrook-ryou}, a Chinese remainder theorem argument is used. However, it is slightly more exotic (though not deeper) than the Chinese remainder theorem found in a typical text on ring theory. In general, the additional abstraction required by the proof in the present leads to an overall streamlining of the proof compared to the one in \cite{hambrook-fraser-Rn}. While the argument is more difficult (due to its generality), it is more efficient. 



\section{Algebraic Number Theory}\label{Algebraic Number Theory: Basics section}

In this section, we give the definitions and facts from algebraic number theory necessary for the proof of our main theorem. As a reference for this material, see \cite{Jarvis14}. 
We also prove some lemmas that we will need. 
 

Recall that $K$ is a number field of degree $d$, $\OK$ is its ring of integers, and $B = \cbr{\omega_1,\ldots,\omega_d}$ is an integral basis for $K$. 


Let $\mathcal{E}_K$ denote the set of all injective homomorphisms 
from $K$ into $\mathbb{C}$. There are precisely $d$ such homomorphisms. 
For each $q \in K$, the field norm of $q$ is 
\begin{align}\label{field norm formula}
N(q) 
=
\prod_{\tau \in \mathcal{E}_K} \tau(q).
\end{align}
and the field trace of $q$ is 
\begin{align}\label{field trace formula}
\tr(q) 
=
\sum_{\tau \in \mathcal{E}_K} \tau(q).
\end{align}
The field norm is multiplicative: 
$N(ab)=N(a)N(b)$ for all $a,b \in K$. 
The field trace is additive: 
$\tr(a+b)=\tr(a) + \tr(b)$ for all $a,b \in K$. 



\begin{lemma}\label{C1 bounds lemma}
For the integral basis $B = \cbr{\omega_1,\ldots,\omega_d}$, define 
\begin{align*}%\label{defn of C1}
C_B = \max\cbr{ \textstyle \rbr{\sum_{i=1}^d |\tau(\omega_i)|^2}^{1/2} : \tau \in \mathcal{E}_K }.
\end{align*}
For every $q \in K$, 
\begin{align}\label{C1 norm bound}
|N(q)| \leq C_B^d |q|^d.
\end{align}
\end{lemma}
\begin{proof}
If $q = \sum_{i=1}^d q_i \omega_i \in K$ and $\tau \in \mathcal{E}_K$, then 
$\tau(q) = \sum_{i=1}^d q_i \tau(\omega_i) \in \CC$.  
The Cauchy-Schwarz inequality implies $|\tau(q)| \leq C_B|q|$, so that \eqref{field norm formula} implies \eqref{C1 norm bound}.
\end{proof}



\begin{lemma}
For all $a,b \in K$, $|ab| \lesssim_B |a||b|$.
\end{lemma}
\begin{proof}
Write $a = \sum_{i=1}^{d} a_i \omega_i$, 
$b = \sum_{j=1}^{d} b_j \omega_j$, and $\omega_i \omega_j = 
\sum_{k=1}^{d} c_{ijk} \omega_k$, where $a_i, b_j, c_{ijk} \in \QQ$. 
%Then $ab = \sum_{k}\sum_{i} \sum_{j} a_i b_j c_{ijk} \omega_k$. 
By the Cauchy-Schwarz inequality, 
$$
|ab|^2 
= 
\sum_{k}  \abs{\sum_i \sum_j a_i b_j c_{ijk} }^2
\lesssim_B \rbr{\sum_i |a_i|^2} \rbr{\sum_j |b_j|^2}
= |a|^2|b|^2. 
$$
\end{proof}



An ideal of $\OK$ is an additive subgroup $I$ of $\OK$ such that $rI = \cbr{ra:a \in I} \subseteq I$ for all $r \in \OK$. A fractional ideal of $\OK$ is a subset $J \subseteq K$ of the form 
$$J = \frac{1}{r}I = \cbr{\dfrac{1}{r}a : a \in I},$$
where $I$ is an ideal of $\OK$ and $r$ is a non-zero element of $\OK$  (neither $I$ nor $r$ are uniquely determined by $J$).  
Every ideal of $\OK$ is a fractional ideal of $\OK$, and every fractional ideal of $\OK$ is an $\OK$-submodule of $K$. 
A fractional ideal of $\OK$ is an ideal of $\OK$ if and only if it is contained in $\OK$. The principal fractional ideal generated by the element $q \in K$ is $\abr{q} = q\OK$. 

The product of two fractional ideals $I$ and $J$ is the fractional ideal
$$
IJ = \cbr{\sum_{i=1}^{n} a_i b_i : a_i \in I, b_i \in J, n \in \NN}. 
$$
The non-zero fractional ideals of $\OK$ form a commutative group under multiplication.  The multiplicative identity is $\OK$. The inverse of a non-zero fractional ideal $J$ is 
$$
J^{-1} = \cbr{a \in K : aJ \subseteq \OK }. 
$$
The group is partially ordered by inclusion
: $I \subseteq J$ implies $HI \subseteq HJ$ for all non-zero fractional ideals $H,I,J$.  
Consequently, if $I \subseteq \OK \subseteq J$, then $J^{-1} \subseteq \OK \subseteq I^{-1}$. 


A non-zero proper ideal $P$ of $\OK$ is called prime if 
$ab \in \OK \setminus P$ whenever $a,b \in \OK \setminus P$. 
It is not generally true that every non-zero, non-unit element of $\OK$ can be uniquely factored into a product of prime elements. 
However, every non-zero proper ideal of $\OK$ can be uniquely factored 
as a product of prime ideals. 



If $J = \frac{1}{r}I$ is a fractional ideal of $\OK$, where $I$ is an ideal of $\OK$ and $r$ is a non-zero element of $\OK$, the ideal norm of $J$ is defined to be 
\begin{align*}%\label{ideal norm fractional} 
N(J) = \dfrac{|\OK/I|}{|\OK/\abr{r}|},    
\end{align*}
which is finite whenever $J$ is non-zero. 
It can be shown that this definition does not depend on the choice of $I$ and $r$. Moreover, the ideal norm is multiplicative: 
For all fractional ideals $J_1,J_2$ of $\OK$, 
$$N(J_1 J_2) = N(J_1)N(J_2).$$ 


\begin{lemma}\label{idealquotient}
If $I_2 \subseteq I_1$ are fractional ideals of $\OK$,  
then  $N(I_2) = N(I_1)|I_1/I_2|$. 
\end{lemma}
\begin{proof}
Choose $r \in \OK \setminus \cbr{0}$ such that $rI_2 \subseteq rI_1 \subseteq \OK$. 
Then 
\begin{align*}
N(\abr{r}) N(I_2)
&=
N(r I_2)
=
|\OK/rI_2|
=
|\OK/rI_1| \cdot |rI_1/rI_2| \\
&=
N(r I_1) |rI_1/rI_2|
=
N(r I_1) |I_1/I_2| 
=
N(\abr{r}) N(I_1) |I_1/I_2|. 
\end{align*}
\end{proof}

The field norm and ideal norm are related as follows:
For all $q \in K$, 
\begin{align}\label{field-norm-ideal-norm-formula}
|N(q)| = N(\abr{q}).  
\end{align}
%
This fact combines with Lemmas \ref{C1 bounds lemma} and \ref{idealquotient} to give the following separation lemma. 

\begin{lemma}\label{separation lemma}
Let $I_1$ and $I_2$ be non-zero fractional ideals of $\OK$. Let $r_1 \in I_1^{-1}$ and $r_2 \in I_2^{-1}$ with $r_1 \neq r_2$. Then 
$$
C_B|r_1 - r_2| 
\geq (N(I_1 \cap I_2))^{-1/d}
\geq (N(I_1) N(I_2))^{-1/d}. 
$$
\end{lemma}

\begin{proof}
Since $I_1 I_2 \subseteq I_1 \cap I_2$, Lemma \ref{idealquotient} implies $N(I_1 I_2) \geq N(I_1 \cap I_2)$, which gives the second inequality. Now we prove the first inequality.
We have $r_1 (I_1 \cap I_2) \subseteq r_1 I_1 \subseteq \OK$ 
and 
$r_2 (I_1 \cap I_2) \subseteq r_2 I_2 \subseteq \OK$. 
Thus $\abr{r_1 - r_2}(I_1 \cap I_2) \subseteq \OK$. Now Lemma \ref{idealquotient} implies 
$N(\abr{r_1 - r_2})N(I_1 \cap I_2) \geq N(\OK) = 1$. 
The desired result follows from Lemma \ref{C1 bounds lemma} and \eqref{field-norm-ideal-norm-formula}.
\end{proof}


We now define the different ideal $\delta$ of $\OK$.  
The set 
$$
\{ x \in K : \tr(xy) \in \ZZ \ \text{ for all } y \in \OK \}
$$ 
is a fractional ideal of $\OK$. 
The different ideal $\delta$ is the \textit{inverse} of this fractional ideal. In other words, $\delta^{-1}$ is equal to the fractional ideal above. Note that $\delta \subseteq \OK \subseteq \delta^{-1}$. 





Recall that the integral basis is $B = \cbr{\omega_1,\ldots,\omega_d}$. % 
Define the $d \times d$ matrix 
$$T =
%[T_{ij}] = 
[\tr(\omega_i \omega_j)].$$ 
Since $(x,y) \mapsto \tr(xy)$ is a non-degenerate bilinear form on $\QQ^d \isom K$, $T$ is invertible and 
the dot product is related to the trace by 
$$\tr(r(T^{-1}s)) = r \cdot s \text{ for all $r,s \in \QQ^d \isom K$}.$$
It follows that $T^{-1}(\ZZ^d) = \delta^{-1}$.

For $u \in \RR$, define 
$$
e(u)=e^{-2\pi i u}. 
$$

\begin{lemma}\label{NEWGeomSeriesProp}
Let $I$ be a non-zero ideal in $\OK$ and let $R(I) \subseteq I^{-1}$ be a complete set of representatives of $ I^{-1} / \OK$. Let $s \in \OK \isom \ZZ^d$. Then 
\begin{align*}
\sum_{r \in R(I)} e(r \cdot s)
=
\left\{
\begin{array}{ll}
N(I) & \text{if $s \in T(\delta^{-1} I) $}  \\
0       & \text{if $s \notin T(\delta^{-1} I) $.} 
\end{array}
\right.
\end{align*}
\end{lemma}


\begin{proof}
For any fractional ideal $J$, we have $J=\OK J$, and hence (as is easily verified)  $J \subseteq \delta^{-1}$ if and only if $\tr(z) \in \ZZ$ for all $z \in J$. 
Taking 
$J = (T^{-1} s) I^{-1}$, it follows that $s \in T(\delta^{-1}I)$ if and only if $r \cdot s =  \tr(r (T^{-1} s)) \in \ZZ$ for all $r \in I^{-1}$. 


Therefore, if $s \in T(\delta^{-1} I)$, we have  
\[
\sum_{r \in R(I)} e(r \cdot s) = \sum_{r \in R(I)} 1 = |R(I)|.
\]
But, by Lemma \ref{idealquotient}, 
\[
|R(I)| = |I^{-1} / \OK| 
= \dfrac{N(\OK)}{N(I^{-1})} 
%= \dfrac{N(I)N(I^{-1}) }{ N(I^{-1}) }
= N(I).
\]

Now assume $s \notin T(\delta^{-1} I)$. 
Then there exists an $r_0 \in I^{-1}$ such that $r_0 \cdot s = \tr(r_0(T^{-1}s)) \notin \ZZ$, and hence $e(r_0 \cdot s) \neq 1$. 
By pigeonholing, $r_0 + R(I)$ is also a complete representatives of $I^{-1}/\OK$.
So there is a bijection $f:R(I) \to r_0 + R(I)$ such that $r - f(r) \in \OK \isom \ZZ^d$ for all $r \in R(I)$. 
Since $r \mapsto e(r \cdot s)$ is $\ZZ^d$-periodic, 
\[
\begin{split}
\sum_{r \in R(I)} e(r \cdot s)  
= \sum_{r \in R(I)} e(f(r) \cdot s)
= e(r_0 \cdot s) \sum_{r \in R(I)} e(r \cdot s).
\end{split}
\]
Since $e(r_0 \cdot s) \neq 1$, 
the sum must be zero.
\end{proof}



We shall need the following version of the Chinese remainder theorem. 

\begin{lemma}\label{concrete crt lemma}
Suppose $I_1$, $I_2$, and $D$ are fractional ideals of $\OK$. 
Suppose $a_1, a_2 \in D$. 
\begin{enumerate}[(a)]
\item If $a_1 - a_2 \notin D(I_1 + I_2)$, then
\[(a_1 + DI_1) \cap (a_2 + DI_2) = \emptyset.\]
\item If $a_1 - a_2 \in D(I_1 + I_2)$, then there exists $a \in D$ such that 
$$
(a_1 + DI_1) \cap (a_2 + DI_2) = a + (D I_1 \cap DI_2).
$$
\end{enumerate}
\end{lemma}

\begin{proof}
\mbox{}\\
\begin{enumerate}[(a)]
%\textbf{Proof of (a)}
    \item 
If $(a_1 + D I_1) \cap (a_2 + D I_2) \neq \emptyset$, then $a_1 + d_1 = a_2 + d_2$ for some $d_1 \in DI_1$ and $d_2 \in DI_2$, and so 
$$
a_1 - a_2 = d_2 - d_1 \in DI_1 + DI_2 = D(I_1 + I_2). 
$$
\item 
%\textbf{Proof of (b)} 
Note 
$$
(I_1 + I_2)^{-1}I_1 + (I_1 + I_2)^{-1}I_2 = (I_1 + I_2)^{-1} (I_1 + I_2) = \OK.$$ 
So there exist $x_1 \in (I_1 + I_2)^{-1}I_1$ and $x_2 \in (I_1 + I_2)^{-1}I_2$ such that $x_1+x_2=1$. 
Define $a = x_1a_2 + x_2a_1$. 
Suppose $a_1 - a_2 \in D(I_1 + I_2)$. 
Then 
\begin{align*}
a &= x_1a_2 + (1-x_1)a_1 = a_1 + (a_2 - a_1)x_1 
\\
&\in
a_1 + D(I_1 + I_2)(I_1 + I_2)^{-1}I_1
=
a_1 + DI_1. 
\end{align*}
A similar calculation shows that $a \in a_2 + DI_2$. The desired equality is now easily verified.
\end{enumerate}
\end{proof}





\section{Intermediate  Functions}\label{sec_Single-
Scale}

In this section, we define several functions necessary for the construction of the measure $\mu$. We also prove necessary Fourier-analytic estimates on these functions. 





\subsection{The Functions \texorpdfstring{$\phi$}{phi} and \texorpdfstring{$\Phi$}{Phi}}


Fix $\phi: \RR^d \to \RR$ such that $\phi$ is $C^{\infty}$, $\phi \geq 0$, $\supp(\phi) \subseteq [-1,1]^d$, and 
$\widehat{\phi}(0) = \int \phi(x) dx = 1$. For each non-zero ideal $I$ of $\OK$ and each $\eta > 0$, define 
\begin{align}\label{Phi defn}
\Phi_{I,\eta}(x) = \sum_{r \in I^{-1}} \eta^{-d} \phi((x-r)/\eta). 
\end{align}
The sum converges uniformly on bounded subsets of $\RR^d$. 
In fact, since $\phi$ is compactly supported, for any bounded subset of $\RR^d$, only finitely many terms of the sum are non-zero on that subset. 
Note that $\eta^{-d} \phi ((x-r)/\eta)$ is an $L^1$-normalized bump function on the 
$\ell^{\infty}$-ball with radius $\eta$ and center $r$. 
Note also that $\Phi_{I,\eta}$ is $C^{\infty}$, non-negative, and $\ZZ^d$-periodic. 



\begin{lemma}\label{Phi-Fourier-Lemma}
Let $I$ be a non-zero ideal of $\OK$, let $s \in \ZZ^d \isom \OK$, and let 
$\eta > 0$. Then 
\begin{equation}\label{Phi-Fourier-eqn-1}
\widehat{\Phi_{I,\eta}}(s) = \widehat{\phi}(\eta s) N(I)\mathbf{1}_{\delta^{-1}I} (T^{-1} s). 
\end{equation}
Moreover, if $0 < \eta < C_B^{-1}N(I)^{-1/d}$, then  
\begin{equation}\label{Phi-Fourier-eqn-2}
\widehat{\Phi_{I,\eta}^2}(s) = \eta^{-d} \widehat{\phi^2}(\eta s) N(I)\mathbf{1}_{\delta^{-1}I} (T^{-1} s).
\end{equation}
\end{lemma}
\begin{proof}
For each prime ideal $I$, set $R(I) = I^{-1} \cap [0,1)^d$. Note that $R(I)$ is a complete set of representatives of $I^{-1} / \OK$. 
Every element $r \in I^{-1}$ can be written uniquely as $r=r'+k$, where $r' \in R(I)$ and $k \in \OK$. Moreover, $\OK$ is identified with $\ZZ^d$ via the integral basis $B$. Thus 
\begin{align*}
\widehat{\Phi_{I,\eta}}(s)
&=
\int_{[0,1)^d} \sum_{r \in I^{-1}}  \eta^{-d} \phi ((x-r)/\eta) e(s \cdot x)  dx
\\
&=
\sum_{r \in I^{-1}}  \int_{[0,1)^d}  \eta^{-d} \phi ((x-r)/\eta) e(s \cdot x)  dx
\\
&=
\sum_{r \in R(I)} \sum_{k \in \ZZ^d}  \int_{[0,1)^d}  \eta^{-d} \phi ((x - r - k)/\eta) e(s \cdot x)  dx
\\
&=
\sum_{r \in R(I)} \int_{\RR^d}  \eta^{-d} \phi ((x - r)/\eta) e(s \cdot x)  dx
\\
&=
\sum_{r \in R(I)}  e(s \cdot r) \int_{\RR^d}  \phi (u) e(\eta s \cdot u)  du
\\ 
&= 
\widehat{\phi} (\eta s) \sum_{r \in R(I)} e(s \cdot r).
\end{align*}
Now Lemma \ref{NEWGeomSeriesProp} implies \eqref{Phi-Fourier-eqn-1}. 
%
%
If $0 < \eta < C_B^{-1}N(I)^{-1/d}$, then Lemma \ref{separation lemma} implies $\phi((x-r)/\eta)$ and $\phi((x-r')/\eta)$ have disjoint support whenever $r$ and $r'$ are distinct elements of $I^{-1}$, and so 
$$
\Phi_{I,\eta}^2(x) = \sum_{r \in I^{-1}} \sum_{r' \in I^{-1}} \eta^{-2d} \phi((x-r)/\eta) \phi((x-r')/\eta) = \sum_{r \in I^{-1}} \eta^{-2d} \phi^2((x-r)/\eta). 
$$ 
Thus an argument similar to the one above gives \eqref{Phi-Fourier-eqn-2}.
\end{proof}


\subsection{The collection of ideals \texorpdfstring{$Q(M)$}{Q(M)}}


For each real number $M > 1$, 
define $Q(M)$ to be the set of prime ideals $I$ of $\OK$ such that $M^d/2 \leq |N(I)| \leq M^d$. 
By the Landau prime ideal theorem, 
\begin{equation}\label{QMsize}
|Q(M)| \sim \frac{M^d}{\log(M)}.
\end{equation}


\begin{lemma}\label{primedivisorbound}
Let $J$ be a non-zero integral ideal of $\OK$. Then 
\begin{equation}\label{primedivisorineq}
|\cbr{I \in Q(M): J \subseteq I}| \leq \frac{\log (2N(J))}{d \log M}.
\end{equation}
\end{lemma}
\begin{proof}
Since $J$ factors uniquely as a product of prime ideals $I$ that contain $J$, and since the ideal norm is multiplicative, 
\[\prod_{\substack{I \in Q(M) \\ J \subseteq I}} N(I) \leq 
N(J).\] 
The desired result follows by taking logarithms and using that each ideal $I \in Q(M)$ satisfies $d\log M \leq \log (2N(I))$. 
\end{proof}







\subsection{The Functions \texorpdfstring{$F_k$}{Fk}}

Let $\{M_k\}_{k=0}^{\infty}$ 
be an increasing sequence of positive real numbers. We assume $M_0$ is large 
and the sequence $\{M_k\}_{k=0}^{\infty}$ increases rapidly. Sufficient size and growth conditions will be noted as they arise in the argument. 

Fix $-d < \rho < d$. For each odd $k$, fix an ideal $J_k \in Q(M_k^{1+\frac{\rho}{d}})$. Thus $J_k$ is a prime ideal  
such that 
$$M_k^{d+\rho}/2 \leq |N(J_k)| \leq M_k^{d+\rho}.$$ 
Assuming $M_1$ is large enough, \eqref{QMsize} guarantees such an ideal exists. Note $J_k \in Q(M_k)$ if $\rho = 0$. 
%

Recall that $\rho^{+}=\max\cbr{\rho,0}$ and 
$\rho^{-}=\max\cbr{-\rho,0}$. 
Define 
$$
\P_k = 
\left\{
\begin{array}{ll}
%1+M_k^{-\rho} 
M_k^{\rho^-}
& \text{if } k \text{ is odd} \\
0 & \text{if } k \text{ is even} \\
\end{array}
\right.
$$
%


Recall that  $\tau > 1$. Define $\eta_k  = M_k^{-(1+\tau)}$. 
%
%
Define $c_k$ by 
$$
\dfrac{1}{c_k}  = \P_k N(J) + \sum_{I \in Q(M_k)}  N(I).
$$
For future reference, by \eqref{QMsize} and assuming $M_k$ is large enough, we have 
\begin{align}\label{ck size}
c_k \sim \log(M_k) M_k^{-2d}. 
\end{align}
%
Define $F_k$ by 
\begin{align}\label{Fk defn}
F_k(x) 
= c_k 
\rbr{ 
\P_k \Phi_{J_k,\eta_k}(x) + \sum_{I \in Q(M_k)} \Phi_{I,\eta_k}(x)  
}. 
\end{align}



\begin{lemma}\label{Fk fourier lemma}
Define 
$$
C_0 = 2^{-1/d}C_B^{-1} \|T^{-1}\|^{-1} N(\delta)^{1/d}. 
$$
Let $k \geq 1$ and $s \in \ZZ^d$. 
%%
%%
\begin{align}
\label{F-1*}
\widehat{F_k}(0) &= 1.  
\\
\label{F-2*}
|\widehat{F_k}(s)| &\leq 1, \qquad \forall s \in \ZZ^d 
\\
\label{F-3*}
|\widehat{F_k}(s)| &=0, \qquad \text{if} \  0 < |s| \leq 
C_0 M_k^{1-\rho^{-}/d}.
\\
\label{F-4*}
|\widehat{F_k}(s)|
&\lesssim M_k^{-d+\rho^+} 
\log|s| \qquad \text{if} \ |s| 
\geq 
C_0 M_k^{1-\rho^{-}/d}.
\\
\label{F-5*}
|\widehat{F_k}(s)|
&\lesssim_N  
M_k^{-d+\rho^+}
\log |s| \left(\frac{M_k^{1+\tau}}{|s|} \right)^N \qquad \text{if} \ |s| \geq M_k^{1+\tau}. 
\end{align}
\end{lemma}

\begin{proof} 
\
By Lemma \ref{Phi-Fourier-Lemma}, we obtain 
$$
\widehat{F_k}(s) 
= c_k \widehat{\phi}(\eta_k s) \left( P_k N(J_k) \mathbf{1}_{\delta^{-1}J_k} (T^{-1} s)+ \sum_{I \in Q(M_k)} N(I) \one_{\delta^{-1} I}(T^{-1}s) \right).
$$

Since $\widehat{\phi}(0) = 1$, \eqref{F-1*} follows immediately. Since $F_k \geq 0$, \eqref{F-2*} follows from \eqref{F-1*}.

We prove the contrapositive of \eqref{F-3*}. 
Assume $s \neq 0$ and $\widehat{F_k}(s) \neq 0$. 
So there is at least one $I \in \cbr{J_{k}} \cup Q(M_k)$ such that $T^{-1} s \in \delta^{-1} I$. 
Then, by Lemma \ref{C1 bounds lemma} and Lemma \ref{idealquotient},
$$
N(\delta^{-1})N(I) \leq N(\abr{T^{-1} s}) = |N(T^{-1} s)| \leq C_B^d |T^{-1} s|^d \leq C_B^d \|T^{-1}\|^d |s|^d. 
$$
If $I \in Q(M_k)$, then $N(I) \geq M_k^d/2$. If $I = J_k$, then $N(I) \geq M_k^{d+\rho}/2$. 
Thus $$|s| \geq 2^{-1/d}C_B^{-1} \|T^{-1}\|^{-1} N(\delta)^{1/d} \min\cbr{M_k,M_k^{1+\rho/d}}.$$  

Now we prove \eqref{F-4*} and \eqref{F-5*}. 
Recall $N(J_k) \leq M_k^{d+\rho}$ and $N(I) \leq M_k^d$ for all $I \in Q(M_k)$. 
By \eqref{ck size}, $c_k \lesssim \log(M_k) M_k^{-2d}$. 
We need an upper bound on the number of $I \in Q(M_k)$ such that $T^{-1}s \in \delta^{-1} I$. 
Note that $T^{-1}s \in \delta^{-1} I$ implies $ (T^{-1}s) \delta \subseteq I$. So it suffices to find an upper bound on the number of $I \in Q(M_k)$ such that $ (T^{-1}s) \delta \subseteq  I$. By applying Lemma \ref{primedivisorbound} with $J = (T^{-1}s) \delta$ and then using \eqref{C1 norm bound} and \eqref{field-norm-ideal-norm-formula}, the number in question is seen to be $\lesssim \log |s|/\log M_k$, assuming $M_k$ (and hence $|s|$) is large enough.  
Therefore,
\[
|\widehat{F_k}(s)| 
\lesssim M_k^{-d} 
( 
%(1+ M_k^\rho)
M_k^{\rho^+}
\log M_k
+
\log |s|) |\widehat{\phi} (\eta_k s)|.
\]
Since $\phi$ is Schwartz, \eqref{F-4*} and \eqref{F-5*} follow immediately. 
\end{proof}

\section{Intermediate Measures}
\label{Intermediate Measures}




In this section, we define and establish Fourier-analytic estimates on measures $\mu_k$ and $\mu_{l,k}$.
These measures are used in the construction of $\mu$ and in proving the required properties of $\mu$.

\subsection{The measures \texorpdfstring{$\mu_k$}{muk}}


Let $\phi_0:\RR^d \to \RR$ be a non-negative Schwartz function with support contained in  $([-3/8,-1/8] \cap [1/8,3/8])^d$ and $\widehat{\phi_0}(0) = \int \phi_0(x) dx = 1$. 
%


We define measures $\mu_k$ on $\RR^d$ by
$$
d\mu_0 = \phi_0 dx, \quad d\mu_k = \phi_0  F_{M_1} \cdots F_{M_{k}} dx \quad \text{ for all } k \in \NN.
$$
To simplify the statement and proof of the next lemma, we define $\mu_{-1}=0$. 



\begin{lemma}\label{mu_k_lemma}
Let $N$ be an integer with $N \geq 100 d$. 
If $(M_{k})_{k=0}^{\infty}$  grows sufficiently rapidly, then for all integers $k \geq 0$ and all $s \in \ZZ^d$, the following statements hold: 
\begin{equation}
\label{mu-1} |\widehat{\mu}_k(s) | \lesssim 1,   
\end{equation}
%
\begin{equation}
\label{mu-2} |\widehat{\mu}_k(s)  -\widehat{\mu}_{k-1}(s)| \lesssim 
M_k^{-(N-2d)}\log(M_k)  \qquad {\text{if}} \ |s| \leq 
% 
\frac{1}{2}C_0 M_k^{1-\rho^{-}/d}, 
\end{equation}
\begin{equation}
\label{mu-3'}
|\widehat{\mu}_k(s)| \lesssim 
%
M_k^{-d+\rho^+}
\log |s|  \qquad {\text{if}} \ |s| \geq 
% 
\dfrac{1}{2} C_0 M_k^{1-\rho^{-}/d}, 
\end{equation}
\begin{equation}
\label{mu-4'}
|\widehat{\mu}_k(s)| \lesssim  
%
M_k^{-d+\rho^+}  
\log |s| \left( \frac{M_k^{1+\tau}}{|s|} \right)^{N-d} \qquad {\text{if}} \ |s| \geq 2M_k^{1+\tau}.    
\end{equation}


\end{lemma}

\begin{proof}
We prove \eqref{mu-1}, \eqref{mu-2}, \eqref{mu-3'}, and \eqref{mu-4'} by induction. 

\textit{Base Step.}  
When $k=0$, \eqref{mu-1}, \eqref{mu-2}, \eqref{mu-3'}, \eqref{mu-4'} follow from 
our choice that $\phi_0$ is Schwartz and $\widehat{\phi_0}(0)=1$, by taking $M_0$ and the implied constants large enough. 


\textit{Inductive Step.} 
Assume \eqref{mu-1}, \eqref{mu-2}, \eqref{mu-3'}, \eqref{mu-4'} hold with $k=0,\ldots,n$.
We shall prove \eqref{mu-1}, \eqref{mu-2}, \eqref{mu-3'}, \eqref{mu-4'} when $k=n+1$ 
using Lemma \ref{convstab lemma}.
Set $G = \widehat{F_{n+1}}$, 
$H = \widehat{\mu_n}$, 
$L(x) = \log(x)$, 
$D = 
M_{n+1}^{-d+\rho^+}
$, 
$A = C_0 M_{n+1}^{1-\rho^{-}/d} $, 
$B = M_{n+1}^{1+\tau}$. 
Notice that the choice $N \geq 100d$ implies that \eqref{Dcond} holds. 
It also implies that $N>2d+2a$ holds for $a = 1$. 
The inequalities \eqref{F-3*}, \eqref{F-4*}, \eqref{F-5*}  in Lemma \ref{Fk fourier lemma} when $k=n+1$ imply (respectively) that the assumptions \eqref{convcond1}, \eqref{convcond2}, \eqref{convcond3} in Lemma \ref{convstab lemma} hold. 
The inequalities \eqref{F-1*}, \eqref{F-2*} in Lemma \ref{Fk fourier lemma} when $k=n+1$ imply $|G| \leq G(0)=1$. 
Our assumption of \eqref{mu-1} when $k=n$ implies $H$ is bounded. 
Also, our assumption of \eqref{mu-4'} when $k=n$ implies that \eqref{convcond4} in Lemma \ref{convstab lemma} holds with an appropriately large implicit constant. 
Because of our choice of $A$ and $B$ above, 
if $M_{n+1}$ is large enough, then the conclusions \eqref{convlemmares1}, \eqref{convlemmares2}, \eqref{convlemmares3} in Lemma \ref{convstab lemma} hold. 
But \eqref{convlemmares1}, \eqref{convlemmares2}, \eqref{convlemmares3} immediately give 
\eqref{mu-2}, \eqref{mu-3'}, \eqref{mu-4'} (respectively) when $k=n+1$. 
It remains to prove \eqref{mu-1} when $k=n+1$. 
Recall we have assumed 
\eqref{mu-2} holds with $k=0,\ldots,n$, and we have proved \eqref{mu-2} when $k=n+1$. 
Therefore 
\begin{align*}
|\widehat{\mu_{n+1}}(0) - \widehat{\mu_0}(0)|
= 
\sum_{k=1}^{n+1} (\widehat{\mu_{k}}(0) - \widehat{\mu_{k-1}}(0))
\lesssim  \sum_{k=1}^{n+1} 
%%%M_{k}^{-(N-2d-1)}\log(M_{k}).
M_{k}^{-(N-2d)}\log(M_{k}).
\end{align*}
However,  
$\widehat{\mu_0}(0) = \widehat{\phi_0}(0) = 1$.
Furthermore, 
since $\mu_{n+1}$ has non-negative density, 
$|\widehat{\mu_{n+1}}(s)| \leq \widehat{\mu_{n+1}}(0)$. 
Thus, 
if $M_{1},\ldots,M_{n+1}$ are large enough, we get \eqref{mu-1} when $k=n+1$. 
%
\end{proof}




\subsection{The measures \texorpdfstring{$\mu_{l,k}$}{mul,k}}

The following measures will be used to prove regularity and the failure of the restriction estimate for $\mu$.   

For $ l \geq k \geq 1$, define $$\mu_{l,k} 
:=  
F_{M_l} \cdots F_{M_{k+1}} F_{M_{k-1}} \cdots F_{M_1} \phi_0
= F_{M_l} \cdots F_{M_{k+1}} \mu_{k-1}. 
$$
Note $\mu_{k,k} := \mu_{k-1}$. 

We can prove an analog of Lemma \ref{mu_k_lemma} for $\mu_{l,k}$ using a virtually identical argument. However, we only need the following weaker version. We omit the proof for brevity. 





\begin{lemma}\label{mu_k_ell_lemma}
For all $l > k \geq 1$ and $s \in \ZZ^d$, 
\begin{equation}\label{mukleq}
\left| \widehat{\mu_{l,k}} (s)  \right| \lesssim 1,  
\end{equation}
\begin{equation}\label{mukleq2}
\left| \widehat{\mu_{l,k}} (s) - \widehat{\mu_{l-1,k}}(s) \right| \leq \log(M_l)M_l^{-100} \qquad \text{if}\  |s| \leq  
\frac{1}{2}C_0 M_{l}^{1-\rho^{-}/d}.
\end{equation}
\end{lemma}


The following estimate will be used several times. 

\begin{lemma}\label{H mu convolution lemma}
Let $l > k \geq 1$.  
Let $C_k > 0$ and $H_k:\ZZ^d \to \RR$ such that, for all $u \in \ZZ^d$, 
\begin{align}\label{H mu convolution 1}
|H_k(u)| \leq C_k (1+\eta_k|u|)^{-10d}. 
\end{align}
Then, for all $s \in \ZZ^d$ such that $|s| \leq M_k^{1+\tau}$, 
\begin{align*}
\sum_{t \in \ZZ^d} |H_k(s-t)| 
| \widehat{\mu_{l,k}} (t) - \widehat{\mu_{k-1}}(t) |
\lesssim C_k M_{k}^{-80(1+\tau)d}. 
\end{align*}
\end{lemma}
\begin{proof}
Since $|s| \leq M_{k}^{1+\tau}$, 
if $|t| \geq M_k^{10(1+\tau)}$, then 
$|s-t| \geq |t|/2$, and so \eqref{H mu convolution 1} implies 
$$
|H_k(s-t)| \lesssim C_k (1+\eta_k |t|)^{-10d}. 
$$
We combine this with \eqref{mu-1} and \eqref{mukleq} to obtain 
\begin{align*}
&\sum_{|t| \geq M_k^{10(1+\tau)}} 
|H_k(s-t)| |\widehat{\mu_{l,k}} (t) - \widehat{\mu}_{k-1}(t)| 
\\
&
\lesssim \sum_{|t| \geq M_k^{10(1+\tau)}} C_k  \left( \frac{M_k^{1+\tau}}{|t|} \right)^{10d} 
\lesssim C_k M_k^{-80(1+\tau)d}.
\end{align*}
If $|t| \leq M_k^{10(1+\tau)}$,
since $M_{k}^{10(1+\tau)} \leq %M_{k+1}(\rho)/2$,
\frac{1}{2}C_0 M_{k+1}^{1-\rho^{-}/d}$, 
\eqref{mukleq2} implies that
$$
|\widehat{\mu_{l,k}} (t) - \widehat{\mu}_{k-1}(t)| \leq \sum_{j=k+1}^l \log(M_j) M_j^{-100} \lesssim M_{k+1}^{-99}.
$$
Then, since \eqref{H mu convolution 1} implies $|H_k| \lesssim C_k$, 
we obtain
\begin{align*}
&\sum_{|t| \leq M_k^{10(1+\tau)}}  |H_k (s-t)| |\widehat{\mu_{l,k}} (t) - \widehat{\mu}_{k-1}(t)| \\
&\lesssim C_k M_k^{10(1+\tau) d} M_{k+1}^{-99} \lesssim C_k M_k^{-80(1+\tau)d}.
\end{align*}
\end{proof}



\section{Construction and Fourier Analysis of  \texorpdfstring{$\mu$}{mu}}\label{sec_measure}


In this section we construct the measure $\mu$ and prove it satisfies 
the desired support and Fourier decay conditions (namely, conditions (i) and (iii) of Theorem \ref{mainthm3}). 

\subsection{Construction of  \texorpdfstring{$\mu$}{mu}}




\begin{lemma}\label{weakconv}
The sequence $(\mu_k)$ converges weakly to a non-zero finite Borel measure $\mu$. 
\end{lemma}

The proof of Lemma \ref{weakconv} requires the Lévy continuity theorem. \begin{thm}[Lévy Continuity Theorem]\label{Levy}
Let $(\nu_k)$ be a sequence of Borel measures on $\RR^d$ such that $0 < c \leq \nu_k(\RR^d) \leq C < \infty$ for all $k$. 
Suppose $(\widehat{\nu_k}(\xi))$ converges for each $\xi \in \RR^d$ 
and 
suppose $(\nu_k)$ is tight (i.e., for every $\epsilon>0$ there is a compact set $K_{\epsilon} \subseteq \RR^d$ such that $\nu_k(\RR^d \setminus K_{\epsilon}) < \epsilon$ for all $k$). Then there exists a finite non-zero Borel measure $\nu$ on $\RR^d$ such that $\nu_k$ converges to $\nu$ weakly (i.e., in distribution). 
\end{thm}

For a proof of this theorem when $\nu_k$ are probability measures, see  \cite[Section 26]{billingsley} for $d=1$ and \cite[Section 26]{billingsley} for $d \geq 1$. The general case follows by considering $a_k\nu_k$,  where $1/a_k = \nu_k(\RR^d)$. 
Alternatively, for a direct proof of the general case, see \cite[Section 8.8]{bogachev}. 

We also need the following lemma. 

\begin{lemma}\label{convseries}
Let $(\nu_k)$ be a sequence of Borel measures with support contained in $[-\frac{1}{2},\frac{1}{2}]^d$ and with $\sup_k \nu_k(\RR^d) < \infty$. If $(\widehat{\nu_k}(s))$ converges for each $s \in \mathbb{Z}^d$, then $(\widehat{\nu_k}(\xi))$ converges for each $\xi \in \mathbb{R}^d$. 
\end{lemma}
\begin{proof}
Let $\psi$ be a Schwartz function on $\RR^d$. 
For every $\xi \in \mathbb{R}^d$, we have (see \cite[Eq. (110)]{wolff-book}) 
\[\widehat{\psi \nu_k}(\xi) = \sum_{s \in \mathbb{Z}^d} \widehat{\nu_k}(s) \widehat{\psi}(\xi - s).\] 
Because $\widehat{\psi}$ is a Schwartz function and because $\sup_k |\widehat{\nu_k}(\xi)| \leq \sup_k \nu_k(\RR^d) < \infty$ for all $\xi \in \mathbb{R}^d$, the Lebesgue dominated convergence theorem lets us take the limit as $k \to \infty$. Thus, for each  $\xi \in \RR^d$, 
\[\lim_{k \to \infty} \widehat{\psi \nu_k}(\xi) = \lim_{k \to \infty} \sum_{s \in \mathbb{Z}^d} \widehat{\nu_k}(s) \widehat{\psi}(\xi - s) = \sum_{s \in \mathbb{Z}^d} \lim_{k \to \infty} \widehat{\nu_k}(s) \widehat{\psi}(\xi - s) \]
and the limit is finite. If $\psi$ is equal to 1 on $[-\frac{1}{2},\frac{1}{2}]^d$, then $\psi \nu_k = \nu_k$ for all $k$. 
\end{proof}
\begin{proof}[Proof of Lemma \ref{weakconv}] 
By construction, $\supp(\mu_k) \subseteq \supp(\phi_0) \subseteq [-\frac{1}{2},\frac{1}{2}]^d$ for all $k$. This also implies the sequence $(\mu_k)$ is tight. 
By \eqref{mu-2}, whenever $k \geq n$ and 
$|s| \leq C_0 M_k^{1-\rho^{-}/d}$,
we have 
\[\left|\widehat{\mu_{k}}(s) - \widehat{\mu_{n}}(s) \right| \leq \sum_{j=n}^{\infty} M_j^{-100} \log(M_j).\]
In particular, $|\widehat{\mu_k}(0) - \widehat{\mu_0}(0)| \leq 1/2$. 
Since $\widehat{\mu_0}(0)= \widehat{\phi_0}(0) = 1$,  we have $1/2 \leq \mu_k(\RR^d) = \widehat{\mu_k}(0) \leq 3/2$ for all $k$. 
Furthermore, since the sum above goes to $0$ as $n \to \infty$, the sequence $(\mu_k(s))$ is Cauchy and hence convergent for each $s \in \ZZ^d$. 
Lemma \ref{convseries} implies $(\mu_k(\xi))$ converges for each $\xi \in \RR^d$. 
The Lévy continuity theorem now implies the desired result. 
\end{proof}






\subsection{Support of \texorpdfstring{$\mu$}{mu}}

\begin{prop}\label{support of mu lemma}
$\supp(\mu) \subseteq \supp(\phi_0) \cap E(K,B,\tau).$ 
\end{prop}
\begin{proof}
Because $\mu_k = \phi_0 F_{1} \cdots F_{k}$ and because $\mu$ is the weak limit of $\mu_k$,  
\[\supp (\mu) \subseteq \bigcap_{k=1}^{\infty} \supp (\mu_k) =  \supp(\phi_0) \cap \bigcap_{k=1}^{\infty} \supp (F_{k}).\]
%
%
%
For even values of $k$, 
$$\supp (F_k) = \bigcup_{I \in Q(M_k)} \supp(\Phi_{I, \eta_k}).$$
If $I \in Q(M_k)$ and $\Phi_{I, \eta_k}(x) > 0$, then 
$\dist(x,I^{-1}) \leq \eta_k \leq N(I)^{-(\tau+1)/d}.$  
Since $I^{-1}$ is discrete, it follows that 
$$
\supp(\Phi_{I,\eta_k}) \subseteq \cbr{x \in \RR^d : \dist(x,I^{-1}) \leq N(I)^{-(\tau+1)/d}} 
$$
for every $I \in Q(M_k)$. 
By putting everything together and using that the sets $Q(M_k)$ are disjoint, we have 
$$
\bigcap_{k=1}^{\infty} \supp(F_k) \subseteq E(K,B,\tau). 
$$
\end{proof}


\subsection{Fourier Decay of \texorpdfstring{$\mu$}{mu}}


\begin{prop}\label{Prop_mu_decay}
For all  $s \in \RR^d$ with $|s| \geq 2$, 
\begin{equation}\label{mu_decay}
|\widehat{\mu}(s)| \lesssim 
\log(|s|)|s|^{-\frac{d-\rho^+}{1 + \tau}}.
\end{equation}
\end{prop}


\begin{proof}
Since $\supp(\mu) \subseteq [-1/2,1/2]^d$, by a standard argument (see, e.g., \cite[Lemma 1, p.252-253]{kahane-book} or \cite[Lemma 9.A.4, p.69]{wolff-book}), we only need to prove \eqref{mu_decay} for $s \in \ZZ^d$. 
Moreover, it 
suffices to prove  that, for all integers $k \geq 0$, 
\begin{equation}\label{mu_k_decay}
|\widehat{\mu_k}(s) | \lesssim 
\log(|s|)|s|^{-\frac{d-\rho^+}{1 + \tau}} 
\quad \text{for all $s \in \ZZ, |s| \geq 2$}
, 
\end{equation}
where the implicit constant  does not depend on $k$. 


We will show \eqref{mu_k_decay} by induction. For $k = 0$, the estimate \eqref{mu_k_decay} follows 
%by taking a  sufficiently large implied constant, 
because $\mu_0 = \phi_0$ is a Schwartz function. Suppose $k \geq 1$ and suppose the statement \begin{equation}\label{mu_j_decay}
|\widehat{\mu_j}(s) | \lesssim 
\log(|s|)|s|^{-\frac{d-\rho^+}{1 + \tau}} \quad \text{for all $s \in \ZZ^d, |s| \geq 2$} 
\end{equation} holds for each $0 \leq j \leq k-1$ 
with implicit  constant independent of $j$. We consider three cases. 

Case 1: $2 \leq |s| \leq 
\frac{1}{2}C_0 M_k^{1-\rho^{-}/d}$. 
Let $0 \leq j_0 \leq k-1$ be such that 
$
\frac{1}{2}C_0 M_{j_0}^{1-\rho^{-}/d} 
\leq |s| \leq 
\frac{1}{2}C_0 M_{j_0+1}^{1-\rho^{-}/d}$,
taking $j_0 = 0$ if 
$|s| \leq 
\frac{1}{2}C_0 M_{0}^{1-\rho^{-}/d}$. 
By applying the triangle inequality, inequality \eqref{mu-2}, and inequality \eqref{mu_j_decay} with $j = j_0$, we have
\begin{align*}
|\widehat{\mu_k}(s)| & \leq |\widehat{\mu_{j_0}}(s)| + \sum_{j = j_0 + 1}^k |\widehat{\mu_j}(s) - \widehat{\mu_{j-1}}(s)| \\
& \lesssim 
\log(|s|)
|s|^{-\frac{d-\rho^+}{1 + \tau}}
+ \sum_{j=j_0 + 1}^k 
%%%M_j^{-(N - 2d-1)}\log(M_j)
M_j^{-(N - 2d)}\log(M_j)
\\
& \lesssim 
\log(|s|)
|s|^{-\frac{d-\rho^+}{1 + \tau}}. 
\end{align*}

Case 2:  
$\frac{1}{2}C_0 M_{k}^{1-\rho^{-}/d} 
\leq |s| \leq 2M_k^{1+\tau}$.  Here \eqref{mu-3'} immediately implies \eqref{mu_k_decay}. 

Case 3: $|s| \geq 2M_k^{1+\tau}$. By our choice of $N$, we have 
$N-d \geq (d-\rho^+)/(1+\tau)$. Then \eqref{mu-4'} gives 
\begin{align*}
|\widehat{\mu}_k(s)| 
&\lesssim 
%%\log(M_k)
\log |s| 
M_k^{-d+\rho^+}\left( \frac{M_k^{1+\tau}}{|s|} \right)^{N-d}
\\
&\leq 
%%\log(M_k)
\log |s| 
M_k^{-d+\rho^+}\left( \frac{M_k^{1+\tau}}{|s|} \right)^{\frac{d-\rho^+}{1+\tau}}
\\
&=
%%\log(M_k)
\log(|s|)
|s|^{-\frac{d-\rho^+}{1+\tau}}.
\end{align*}
This implies \eqref{mu_k_decay} for $|s| \geq 2M_k^{1+\tau}$. 
\end{proof}



\section{Regularity of \texorpdfstring{$\mu$}{mu}}\label{regularity_section}




In this section, we prove that $\mu$ satisfies (ii) of Theorem \ref{mainthm3}. We need several preparatory lemmas before getting to the main result, Proposition \ref{Prop_mu_reg}. 

\begin{lemma}\label{separation facts}
    Let $k \in \NN$.
\begin{enumerate}[(a)]
    \item Suppose $r_1, r_2 \in J_k^{-1}$. If $r_1 \neq r_2$,  then $|r_1 - r_2| \gtrsim M_{k}^{-1-\frac{\rho}{d}}$. 
    \item 
    %Suppose $(I_1,r_1) \in Q(M_k) \times I_1^{-1}$ and $(I_2,r_2) \in Q(M_k) \times I_2^{-1}$. 
    Suppose $I_1,I_2 \in Q(M_k)$, $r_1 \in I_1^{-1}$, $r_2 \in I_2^{-1}$. 
    If $r_1 \neq r_2$, then $|r_1 - r_2| \gtrsim M_{k}^{-2}$. 
     \item Suppose $I_1,I_2 \in Q(M_k)$, 
     $r_1 \in I_1^{-1}$, $r_2 \in I_2^{-1}$. 
     If $r_1 = r_2$ and $I_1 \neq I_2$, then 
     $r_1=r_2 \in \OK \isom \ZZ^d$ and (consequently) 
     $\dist(r_1,\supp(\phi_0))=
     \dist(r_2,\supp(\phi_0)) > 3\eta_k$. 
\end{enumerate}
\end{lemma}
\begin{proof}
Lemma \ref{separation lemma} directly implies (a) and (b). 
Now we prove (c). 
Since $I_1$ and $I_2$ are distinct prime ideals of $\OK$, they are relatively prime, i.e., $1 \in I_1+I_2$.
Set $r = r_1=r_2$. 
Since $r \in I_1^{-1}$ and $r \in I_2^{-1}$, we have $rI_1 \subseteq \OK$ and $rI_2 \subseteq \OK$. 
Thus 
$
r = r \cdot 1 \in r(I_1+I_2) = rI_1 + rI_2 \subseteq 
\OK \isom \ZZ^d. 
$
Recalling that $\supp(\phi_0) \subseteq ([-3/8,-1/8] \cap [1/8,3/8])^d$ gives the result. 
\end{proof}



\begin{lemma}\label{Fk nonzero terms lemma}
Let $k \in \NN$. Let  $G_k:\RR^d \to \RR$. 
Suppose $\supp(G_k)$ has diameter $\leq 2\eta_k$ and 
suppose $\supp(G_k)$ intersects $\supp(\phi_0)$. 
Then there exists $I \in Q(M_k)$, $r' \in I^{-1}$, and $r'' \in J_k^{-1}$ 
such that, for all $x \in \RR^d$, 
\begin{align*}
F_k(x) G_k(x) = c_k \eta_k^{-d} 
\rbr{ \phi\rbr{\dfrac{x-r'}{\eta_k}} + P_k \phi\rbr{\dfrac{x-r''}{\eta_k}}}
G_k(x). 
\end{align*}
\end{lemma}
%
\begin{proof}
For brevity, write $\phi_{r,\eta_k}(x)=\phi((x-r)/\eta_k)$. 
By \eqref{Phi defn} and \eqref{Fk defn}, 
\begin{align}\label{Fk expanded 3}
F_k G_k = 
\sum_{r \in J_k^{-1}} P_k c_k  \eta_k^{-d} \phi_{r,\eta_k} G_k + 
\sum_{I \in Q(M)} \sum_{r \in I^{-1}} c_k  \eta_k^{-d} \phi_{r,\eta_k} G_k. 
\end{align}
We assume $(M_k)$ has been chosen so that  $\eta_k=M_k^{-(1+\tau)}$ is much smaller than $M_{k}^{-2}$
and $M_k^{-1-\rho/d}$. 
Note $\supp(\phi_{r,\eta_k}) \subseteq B(r,\eta_k)$. 
By Lemma \ref{separation facts}(a), $\supp(G_k)$ intersects $\supp(\phi_{r,\eta_k})$ for at most one $r \in J_k^{-1}$. 
Thus  
the first sum in \eqref{Fk expanded 3} has at most one non-zero term. 
Now consider two distinct pairs $(I_1,r_1)$ and $(I_2,r_2)$, where 
$I_1,I_2 \in Q(M_k)$, $r_1 \in I_1^{-1}$, $r_2 \in I_2^{-1}$. 
If $r_1 \neq r_2$, Lemma \ref{separation facts}(b) implies 
$\supp(G_k)$ intersects at most one of $\supp(\phi_{r_1,\eta_k})$ and $\supp(\phi_{r_2,\eta_k})$. 
If $r_1 = r_2$ and $I_1 \neq I_2$, Lemma \ref{separation facts}(c) implies 
$$
\dist(\supp(\phi_{r_1,\eta_k}),\supp(\phi_0))
= 
\dist(\supp(\phi_{r_2,\eta_k}),\supp(\phi_0))
> 
\eta_k,$$
so $\supp(G_k)$ intersects neither of $\supp(\phi_{r_1,\eta_k})$ and $\supp(\phi_{r_2,\eta_k})$. 
We conclude that the second sum in \eqref{Fk expanded 3} has at most one non-zero term. 
\end{proof}



\begin{lemma}\label{Fphi_decay}
Let $\psi$ be a Schwartz function with compact support. 
Let $k \in \NN$. Let $B = B(x_0,\eta_k)$. Let $\psi_{B}(x) = \psi((x-x_0)/\eta_k)$. 
Suppose $\supp(\psi_B)$ intersects $\supp(\phi_0)$.
For all $Z \geq 0$ and $s \in \ZZ^d$, 
$$
\widehat{F_k \psi_{B}}(s)
\lesssim_{Z,\phi, \psi}c_k (1+P_k)(1+\eta_k|s|)^{-Z}.
$$
We emphasize that the implied constant does not depend on $x_0$ or $k$. 
\end{lemma}
\begin{proof}
In light of Lemma \ref{Fk nonzero terms lemma}, it suffices to estimate 
$\widehat{fg}(s)$, where $f(x) = \phi\rbr{({x-r})/{\eta_k}}$ and $g(x) = \psi\rbr{({x-x_0})/{\eta_k}}$. 
Note
$$
|\widehat{fg}(s)| 
= 
|\widehat{f} \ast \widehat{g}(s)| 
\leq 
\int_{T_1} |\widehat{f}(t)| |\widehat{g}(s-t)| dt
+
\int_{T_2} |\widehat{f}(t)| |\widehat{g}(s-t)| dt, 
$$
where $T_1 = \cbr{t \in \RR^d: |s|/2 \leq |t|}$ and $T_2 = \cbr{t \in \RR^d: |s|/2 \leq |s-t|}$. 
Also observe that 
$|\widehat{f}(\xi)| = \eta_k^d |\widehat{\phi}(\eta_k \xi)|$ and 
$|\widehat{g}(\xi)| = \eta_k^d |\widehat{\psi}(\eta_k \xi)|$. 
To bound the integral over $T_2$, use that 
$|\widehat{g}(s-t)| \lesssim_{Z,\psi} \eta_k^d (1+\eta_k|s-t|)^{-Z} \leq \eta_k^d (1+\eta_k|s|/2)^{-Z}$ for $t \in T_2$ and that $\|\widehat{f}\|_{L^1}=\|\widehat{\phi}\|_{L^1}$. The integral over $T_1$ is treated similarly. 
\end{proof}


\begin{lemma}\label{Fk bound lemma}
Let $k \in \NN$.  Let $x_0 \in \RR
^d$. 
\begin{enumerate}[(a)]
\item 
If $x_0 \in \supp(\phi_0)$, then 
\begin{align*}
F_k(x_0) \leq c_k \eta_k^{-d}(1+P_k) \|\phi\|_{\infty}. 
\end{align*}
\item 
If $x_0 \in \supp(\phi_0)$ and $x_0 \notin \supp(\Phi_{J_k,\eta_k})$, 
then 
$$
F_k(x_0) \leq c_k \eta_k^{-d} \|\phi\|_{\infty}. 
$$
\end{enumerate}
\end{lemma}
\begin{proof}
Take $G_k = \one_{\cbr{x_0}}$ in Lemma \ref{Fk nonzero terms lemma} and note that the support of $\Phi_{J_k,\eta_k}$ is the union of the closed balls $B(r,\eta_k)$ for $r \in J_k^{-1}$. 
\end{proof}


\begin{lemma}\label{k-ball regularity}
Let $k \in \NN$. Let $B$ be any ball of radius $\eta_k$. 
Let $2B$ be the ball with the same center but twice the radius. 
Then 
\begin{equation} \label{mu(B)-3}
\mu(B) \lesssim \log^2(M_k) 
M_k^{-2d+\rho^-}.
\end{equation}
Furthermore, if $2B$ does \textit{not} intersect $\supp(\Phi_{J_k,\eta_k})$, then 
\begin{equation} \label{mu(B)-1}
\mu(B) \lesssim \log^2(M_k) M_k^{-2d}
\end{equation}
and
\begin{equation}\label{mu(B)-2}
\mu(B) \lesssim \log^2(M_{k-1}) 
%M_{k-1}^{d(\tau-1)}(1+M_{k-1}^{-\rho})
M_{k-1}^{d(\tau-1)+\rho^-}
|Q(M_k)|^{-1} M_k^{-d}.
\end{equation}
\end{lemma}
\begin{proof}
Note \eqref{mu(B)-1} follows from \eqref{mu(B)-2}, if $M_k$ is large enough. So we only need to prove \eqref{mu(B)-3} and \eqref{mu(B)-2}. 

\textit{Step 1: Estimate for $\mu_k(2B)$.} 
We first prove \eqref{mu(B)-3} and \eqref{mu(B)-2} with $\mu$ replaced by $\mu_k$ and $B$ replaced by $2B$. 
Note 
\begin{align*}
\mu_k(2B) = \int_{2B} \phi_0 F_1 \cdots F_k dx. 
\end{align*}
Since the Lebesgue measure of $2B$ is $\lesssim M_k^{-(1+\tau)d}$, it will suffice to obtain appropriate pointwise bounds on the integrand.   
%
If $x \in \supp(\phi_0)$, then, by Lemma \ref{Fk bound lemma}(a), 
\begin{align*}
\phi_0 F_1 \cdots F_k(x) 
&\lesssim \log(M_k) c_k \eta_k^{-d}(1+P_k) \\
\nonumber 
&\lesssim \log^2(M_k) 
M_k^{d(\tau-1)+\rho^-}, 
\end{align*}
assuming $M_k$ is large enough. 
This proves \eqref{mu(B)-3} with $\mu$ replaced by $\mu_k$ and $B$ replaced by $2B$.  
%
Now suppose $2B$ does not intersect $\supp(\Phi_{J_k,\eta_k})$. 
For each $x \in \supp(\phi_0) \cap 2B$, we apply Lemma \ref{Fk bound lemma}(a) to $F_1(x), \ldots, F_{k-1}(x)$ and apply Lemma \ref{Fk bound lemma}(b) to $F_k(x)$ to get  
\begin{align*}
\phi_0 F_1 \cdots F_k(x) &\lesssim  \log(M_{k-1}) c_{k-1} \eta_{k-1}^{-d} (1+P_{k-1}) c_{k} \eta_{k}^{-d} \\
\nonumber 
&\lesssim \log^2(M_{k-1}) 
M_{k-1}^{d(\tau-1)+\rho^-}
|Q(M_k)|^{-1}M_k^{\tau d}, 
\end{align*}
assuming $M_{k-1}$ and $M_k$ are large enough. 
This proves \eqref{mu(B)-2} with $\mu$ replaced by $\mu_k$ and $B$ replaced by $2B$. 


\textit{Step 2: Estimate for $\mu(B)$.} 
We assume $B$ intersects the support of $\phi_0$ (otherwise $\mu_l(B)=0$ and there is nothing to prove). 
Suppose $B = B(x_0,\eta_k)$. So $2B = B(x_0,2\eta_k)$. 
Let $\psi$ be a Schwartz function with support contained in $B(0,2)$, with $\psi \geq 0$ everywhere, and with $\psi \gtrsim 1$ on $B(0,1)$.
Let $\psi_{B}(x) = \psi((x-x_0)/\eta_k)$. 
For arbitrary $l \geq k$, we write
$$
\mu_l(\psi_B) : = \int \psi_B \mu_l dx.
$$
By Plancherel's theorem, we have
$$
\mu_l(\psi_B) - \mu_k (\psi_B)= \sum_{t \in \ZZ^d} \widehat{F_k\psi_B} (t) (\widehat{\mu_{l,k}} (t) - \widehat{\mu}_{k-1}(t)). 
$$
Note that hypotheses of Lemma \ref{Fphi_decay} are satisfied. 
By Lemma \ref{Fphi_decay} with $Z=10d$, the hypotheses of 
Lemma \ref{H mu convolution lemma} are satisfied with $C_k = c_k (1+P_k)$, $H(u) = \widehat{F_k \psi_B}(-u)$, and $s=0$. From Lemma \ref{H mu convolution lemma}, we obtain 
$$
\mu_l(\psi_B) - \mu_k(\psi_B) \lesssim c_k (1+P_k) M_k^{-80(1+\tau)d}.
$$
Since $\one_{B} \lesssim \psi_B \lesssim \one_{2B}$, we have $\mu_{l}(B) \lesssim \mu_{l}(\psi_B)$
and 
$\mu_k(\psi_B) \lesssim \mu_k(2B)$. 
Therefore 
$$
\mu_l(B) 
\lesssim \mu_k(2B) + c_k (1+P_k) M_k^{-80(1+\tau)d}.
$$
Letting $l \rightarrow \infty$, we have established 
\eqref{mu(B)-2} and \eqref{mu(B)-3}. 
\end{proof}



Now we prove the regularity of the measure $\mu$. 




\begin{prop}\label{Prop_mu_reg}
For each ball $B$ in $\RR^d$ of radius $0 < r < 1$,
\begin{equation*}
\mu(B) \lesssim 
\log^2(r^{-1})r^{\frac{2d-\rho^-}{1+\tau}}.
\end{equation*}
\end{prop}



\begin{proof}
\textit{Preliminaries.} Since $\mu$ is a probability measure, the result holds for $r \geq M_0^{-(1+\tau)}$ by taking a large enough implied constant. 
Thus we assume that $M_k^{-(1+\tau)} \leq r \leq M_{k-1}^{-(1+\tau)}$ for some $k \geq 1$. 

We introduce some terminology. 
A $k$-ball is closed ball of radius $\eta_k$ with center $r \in I^{-1}$ for some $I \in Q(M_k) \cup \cbr{J_k}$. 
We consider two kinds of $k$-balls. 
Let $\mathcal{J}_2$ be the set of $k$-balls $J$ such that $2J$ intersects $\supp(\Phi_{J_k,\eta_k})$ and let $\mathcal{J}_1$ be the remaining $k$-balls. (Recall that $2J$ is the ball with the same center as $J$ but with twice the radius.) Note that if $k$ is even, then $\mathcal{J}_2$ is empty. 

The support $\mu$ is contained in the support of $F_k$, which is the union of all the $k$-balls. 
Hence, we have the decomposition
$$
\mu(B) \leq \sum_{J \in \mathcal{J}_1, J \cap B \neq \emptyset} \mu(J) +\sum_{J \in \mathcal{J}_2, J \cap B \neq \emptyset} \mu(J).
$$

\textit{Sum over $J \in \mathcal{J}_1.$} 
First we estimate the sum over $J \in \mathcal{J}_1$. 
Since each ball in $\mathcal{J}_1$ is centered at an element of $I^{-1}$ for some $I \in Q(M_k)$, we will use Lemma \ref{separation lemma} to bound the number of terms in the sum. 
We consider three cases. 

Case 1: $M_k^{-(1+\tau)} \leq r \leq M_k^{-2} $. 
By Lemma \ref{separation lemma}, the centers of balls in $\mathcal{J}_1$ are separated by $\gtrsim M_k^{-2}$. 
Thus the number of terms in the sum over $\mathcal{J}_1$ is $\lesssim 1$. By \eqref{mu(B)-1}, 
$$
\sum_{J \in \mathcal{J}_1, J \cap B \neq \emptyset} \mu(J) 
\lesssim \log^2(M_k) M_k^{-2d} \lesssim \log^2(r^{-1}) r^{\frac{2d}{1+\tau}}.
$$

Case 2: $M_k^{-2} \leq r \leq M_k^{-1} $. Again, by Lemma \ref{separation lemma}, the centers of balls in $\mathcal{J}_1$ are separated by $\gtrsim M_k^{-2}$. Now the number of terms in the sum is $\lesssim (M_k^2r)^d$. By \eqref{mu(B)-1}, 
$$
\sum_{J \in \mathcal{J}_1, J \cap B \neq \emptyset} \mu(J) 
\lesssim \log^2(M_k) r^d \lesssim \log^2(r^{-1}) r^d \leq \log^2(r^{-1}) r^{\frac{2d}{1+\tau}}. 
$$

Case 3: $M_k^{-1} \leq r \leq M_{k-1}^{-(1+\tau)}$. Fix an ideal $I \in Q(M_k)$. By Lemma \ref{separation lemma}, the elements in $I^{-1}$ are at least $M_k^{-1}$ separated. 
So the number of $k$-balls in $\mathcal{J}_1$ that intersect $B$ and whose centers belong to $I^{-1}$ 
is $\lesssim (M_k r)^d$. 
Therefore the total number of $k$-balls in $\mathcal{J}_1$ that intersect $B$ is $\lesssim |Q(M_k)| (M_k r)^d$. 
Using  \eqref{mu(B)-2} and that $r \leq M_{k-1}^{-(1+\tau)}$, we get 
\begin{align*}
\sum_{J \in \mathcal{J}_1, J \cap B \neq \emptyset} \mu(J) \lesssim \log^2(M_{k-1}) 
M_{k-1}^{(\tau -1)d+\rho^-} 
r^d
 \lesssim \log^2 (r^{-1}) 
 r^{\frac{2d-\rho^-}{1+\tau}}. 
\end{align*}

\textit{Sum over $J \in \mathcal{J}_2.$} 
Now we consider the sum over $J \in \mathcal{J}_2$. 
Let $\mathcal{J}'_2$ be the subcollection of $\mathcal{J}_2$ consisting of the balls whose center is an element of $J_{k}^{-1}$. 
Let $\mathcal{J}''_2$ be the remaining balls in $\mathcal{J}_2$. 
By Lemma \ref{separation facts}(b), each ball in $\mathcal{J}'_2$ intersects at most one ball in $\mathcal{J}''_2$. 
Thus the cardinality of $\mathcal{J}_2$ is at most twice that of $\mathcal{J}'_2$. 
Moreover, by Lemma \ref{separation facts}(a), 
the balls in $\mathcal{J}'_2$ are separated by $\gtrsim M_k^{-1-\frac{\rho}{d}}$. 
Thus the number of terms in the sum over $\mathcal{J}_2$ is $ \lesssim \max\cbr{ 1, M_k^{d+\rho} r^d}$.




By \eqref{mu(B)-3}, 
\begin{equation}\label{meas_J_2}
\sum_{J \in \mathcal{J}_2, J \cap B \neq \emptyset} \mu(J) \lesssim \log^2(M_k)
%M_k^{-2d} (1+M_k^{-\rho}) 
M_k^{-2d+\rho^-}
\cdot 
\max\cbr{1, M_k^{d+\rho} r^d}.
\end{equation}
If the maximum 
in \eqref{meas_J_2} 
is $1$, then $M_k^{-(1+\tau)} \leq r \leq M_k^{-1-\frac{\rho}{d}}$, and so 
$$
\sum_{J \in \mathcal{J}_2, J \cap B \neq \emptyset} \mu(J) \lesssim \log^2(r^{-1}) 
r^{\frac{2d+\rho^-}{1+\tau}}.
$$
If the maximum 
in \eqref{meas_J_2} 
is $M_k^{d +\rho} r^d$, 
then 
$$
\sum_{J \in \mathcal{J}_2, J \cap B \neq \emptyset} \mu(J) 
\lesssim 
\log^2(M_k) 
M_k^{-d+\rho^+}
r^d. 
$$
Since the function  
$\log^2(x)x^{-d+\rho^+}$ 
is eventually decreasing and 
since $M_k^{-(1+\tau)} \leq r \leq M_{k-1}^{-(1+\tau)}$, 
the last expression is 
\begin{align*}
\lesssim 
\log^2(r^{-1}) 
r^{\frac{d-\rho^+}{1+\tau}} 
r^d 
\leq 
\log^2(r^{-1})  
r^{\frac{2d-\rho^-}{1+\tau}}. 
\end{align*}
\end{proof}





\section{Failure of the Extension Estimate}\label{sec_Failure_Estimate}

In this section, we show that $\mu$ satisfies (iv) of Theorem \ref{mainthm3} with the functions $f_k = \Phi_{J_k, \eta_k}$. 
Explicitly, our goal is to prove the following proposition. 

\begin{prop}\label{res fail prop}
For all $1 \leq q < \infty$ and all $p < p(\tau, \rho, q, d) := \frac{q(d \tau - \rho)}{(q-1) (d-\rho^+)}$, 
$$
\lim_{\substack{k \to \infty \\ k \text{ odd}}} \dfrac{\|\widehat{\Phi_{J_k,\eta_k} \mu}\|_{L^p}}{\|\Phi_{J_k,\eta_k}\|_{L^q(\mu)}} = \infty.
$$
\end{prop}


Our proof of this proposition requires establishing several intermediate results.  

\subsection{Pointwise Lower Bound}\label{Pointwise Lower Bound}



Fix a constant $0 < c < 1$ such that $\widehat{\phi^2}(\xi) \geq c$ for all $|\xi| \leq c$. (Such a constant exists because $\widehat{\phi^2}(0) = \int \phi^2(x) dx > 0$ and $\widehat{\phi^2}$ is continuous.)  
%
Fix an odd positive integer $k$. 
We will prove:  

\begin{lemma}\label{lem_lb_fkmu}
If $T^{-1} s \in \delta^{-1} J_k$ and $|s| \leq c M_k^{1 + \tau}$, then 
\begin{equation}\label{eq_lb_fkmu}
\left|\widehat{\Phi_{J_k, \eta_k} \mu}(s)\right| \gtrsim \log(M_k) 
M_k^{d \tau+\rho^+}.
\end{equation}
\end{lemma}

The rest of the subsection is devoted to the proof of this lemma. It suffices to show that 
\begin{equation}\label{f_kmu_l_lb}
|\widehat{\Phi_{J_k, \eta_k} \mu_l} (s)| \gtrsim \log(M_k) 
%M_k^{d \tau }(1+M_k^\rho) 
M_k^{d \tau+\rho^+} 
\end{equation}
for $l \geq k$. 
%
Fix $l, k \in \NN$ with $l \geq k$. 
%
Write $\Phi_{J_k, \eta_k}F_k = h_1+h_2$ where
$$
h_1(x) = c_k P_k (\Phi_{J_k, \eta_k}(x))^2
$$
and
$$
h_2(x) = c_k \Phi_{J_k, \eta_k}(x) \sum_{I \in Q(M_k)} \Phi_{I, \eta_k}(x).
$$
Recall that 
$\mu_{l,k} = F_l F_{l-1} \cdots F_{k+1} \mu_{k-1}$. 
Therefore  $\Phi_{J_k, \eta_k} \mu_l  = h_1 \mu_{l,k} + h_2 \mu_{l,k}$, and hence 
\begin{align}\label{phi mu decomp}
\widehat{\Phi_{J_k, \eta_k} \mu_l}(s)
=
\widehat{h_1}(s) \widehat{\mu_{l,k}}(0)
+
\sum_{\substack{t \in \ZZ^d \\  t \neq 0}} \widehat{h_1}(s-t) \widehat{\mu_{l,k}}(t)
+
\sum_{t \in \ZZ^d} \widehat{h_2}(s-t) \widehat{\mu_{l,k}}(t). 
\end{align}



\begin{lemma}
If $T^{-1} s \in \delta^{-1} J$ and $|s| \leq c M_k^{1+\tau}$, then 
\begin{equation}\label{h_1_lb}
|\widehat{h_1} (s) \cdot  \widehat{\mu_{l,k}} (0) | \gtrsim \log(M_k) 
M_k^{d \tau+\rho^+}.
\end{equation}
\end{lemma}
\begin{proof}
% 
By \eqref{Phi-Fourier-eqn-2},  
%
\eqref{ck size}, and the definition of $c$, 
\begin{align*}
|\widehat{h_1}(s)| 
= c_k P_k \eta_k^{-d}N(J_k) |\widehat{\phi^2}(\eta_k s)| 
\gtrsim  
\log(M_k) 
M_k^{d\tau+\rho^+}.
\end{align*}
%
%
Since $\widehat{\mu_{0}}(0) \widehat{\phi}_0(0)=1$, we iterate \eqref{mukleq2} and obtain  $|\widehat{{\mu}_{l,k}}(0)| \gtrsim 1$. 
\end{proof}


\begin{lemma}
If $T^{-1} s \in \delta^{-1} J$ and $|s| \leq M_k^{1+\tau}$, then 
\begin{equation}\label{h_1_ub}
\sum_{t \neq 0 }|\widehat{h_1} (s-t)   \widehat{\mu_{l,k}} (t) | \lesssim \log(M_k) 
M_k^{d \tau + \rho^+ - 80(d+\rho)}.
\end{equation}
\end{lemma}


\begin{proof}
By \eqref{Phi-Fourier-eqn-2} and \eqref{ck size}, 
\begin{align*}
\widehat{h_1}(u) & = c_k P_k \widehat{\Phi_{J_k, \eta_k}^2}(u) \\
& \lesssim c_k \eta_k^{-d} P_k \widehat{\phi^2}(\eta_k u) N(J_k) \mathbf{1}_{\delta^{-1} J_k}  (T^{-1} u) \\
& \lesssim 
\log (M_k) 
M_k^{d \tau+\rho^+} 
\left(1 + \frac{|u|}{M_k^{1 + \tau}} \right)^{-10d} \mathbf{1}_{\delta^{-1} J_k}(T^{-1} u).
\end{align*}
%
Therefore the hypotheses of Lemma \ref{H mu convolution lemma} are satisfied with 
$C_k = \log (M_k) 
%M_k^{d \tau}  (1 + M_k^{\rho})
M_k^{d \tau+\rho^+}$  
and $H_k(u) = \widehat{h_1}(u) \one_{\ZZ^d \setminus \cbr{s}}(u)$.  
Consequently, using the assumption $|s| \leq M_k^{1+\tau}$, we get 
$$
\sum_{t \neq 0} |\widehat{h_1}(s-t) \widehat{\mu_{l,k}}(t)|
\lesssim C_k M_{k}^{-80(1+\tau)d} + \sum_{t \neq 0} |\widehat{h_1}(s-t) \widehat{\mu_{k-1}}(t)|. 
$$
Since $T^{-1}s \in \delta^{-1} J$ and since $\widehat{h}_1(s-t)$ is zero when $T^{-1}(s-t) \notin \delta^{-1} J$, we conclude that $T^{-1}t \in \delta^{-1} J$ for those $t$ such that $\widehat{h_1}(s - t) \neq 0$. 
Thus, by \eqref{C1 norm bound} and Lemma \ref{idealquotient}, 
$|t|^d \gtrsim N(J) \gtrsim M_k^{d+ \rho}$, 
which implies $|t| \geq 2M_{k-1}^{1+\tau}$. 
Thus, by \eqref{mu-4'}, we have
\begin{align*}
|\widehat{\mu}_{k-1}(t) | 
\lesssim 
%\log(M_{k-1}) 
M_{k-1}^{-d+\rho^+} 
\log(|t|)  ( M_{k-1}^{1+\tau} / |t|)^{99d}. 
\end{align*}
Recall that $|\widehat{h}_1 (u)| \lesssim \log(M_k) 
M_k^{d \tau + \rho^+} 
= C_k$. 
Therefore 
$$
\sum_{t \neq 0} |\widehat{h_1}(s-t) \widehat{\mu_{k-1}}(t)|
= 
\sum_{|t| \gtrsim M_k^{1+\frac{\rho}{d}}} |\widehat{h_1}(s-t) \widehat{\mu_{k-1}}(t)|
\lesssim C_k M_k^{-97(d + \rho)}. 
$$
Since $80(1+\tau)d$ and $97(d+\rho)$ are both larger than $80(d+\rho)$, 
we have established \eqref{h_1_ub}. 
\end{proof}



To bound the sum involving $\widehat{h_2}$ in \eqref{phi mu decomp}, we first prove the following lemma. 
%
\begin{lemma}\label{h_2_est}
Suppose $\rho \neq 0$. 
For all 
$s \in \ZZ^d$, 
$$
|\widehat{h_2}(s)| \lesssim  M_k^{d\tau} (1+ M_k^{\rho -d(\tau-1)})(1 + \eta_k |s|)^{-10d}.
$$
\end{lemma}
\begin{proof}
Note that
\begin{align*}
\widehat{h_2}(s) 
= \sum_{I \in Q(M_k)} \sum_{t \in \mathbb{Z}^d} 
c_k \widehat{\Phi_{I, \eta_k}}(t) \widehat{\Phi_{J_k, \eta_k}}(s - t).
\end{align*}
By Lemma \ref{Phi-Fourier-Lemma}, 
\[\widehat{h_2}(s) = \sum_{I \in Q(M_k)} \sum_{t \in \mathcal{T}(I)} 
c_k 
N(I) 
 N(J_k) 
 \widehat \phi(\eta_k t)
\widehat \phi(\eta_k(s - t)),\]
where 
$$
\mathcal{T}(I)
=
\cbr{t \in \ZZ^d : T^{-1}t \in \delta^{-1}I \cap (T^{-1}s + \delta^{-1}J_k)}.
$$
Since at least one of $|s - t|$ and $|t|$ will be $\geq |s|/2$, 
the Schwartz bounds on $\widehat{\phi}$ 
imply 
$$
|\widehat \phi(\eta_k t)
\widehat \phi(\eta_k(s - t))| 
\lesssim (1+\eta_k|s|)^{-10d}(1+\eta_k|t|)^{-50d}.
$$
Therefore,
\begin{equation}\label{h2hatsum2}
|\widehat{h_2}(s)| \lesssim 
(1 + \eta_k |s|)^{-10d} \log (M_k) M_k^{\rho} 
\sum_{I \in Q(M_k)} 
\sum_{t \in \mathcal{T}(I)} 
(1+\eta_k|t|)^{-50d}. 
\end{equation}
Consider a fixed $I \in Q(M_k)$. 
By Lemma \ref{concrete crt lemma},  either the set 
$\delta^{-1} I \cap (T^{-1} s + \delta^{-1} J_k)$ is empty or 
there exists $a \in \delta^{-1}$ such that 
$$
\delta^{-1} I \cap (T^{-1} s + \delta^{-1} J_k)
=
a + \delta^{-1} I \cap \delta^{-1}J_k. 
$$ 
In the first case, $\mathcal{T}(I)$ is empty, so the sum over $t \in \mathcal{T}(I)$ in \eqref{h2hatsum2} is zero. 
In the second case, 
$\delta^{-1} I \cap (T^{-1} s + \delta^{-1} J_k)$ is a translate of 
$\delta^{-1} I \cap \delta^{-1}J_k$, 
so the separation between the elements of $\delta^{-1} I \cap (T^{-1} s + \delta^{-1} J_k)$
 is the same as the separation between the elements of $\delta^{-1} I \cap \delta^{-1}J_k$. 
By Lemma \ref{separation lemma} (with $I_1 = I_2 = ( \delta^{-1} I \cap \delta^{-1}J_k )^{-1}$), the elements of $\delta^{-1} I \cap \delta^{-1}J_k$ are separated by 
$
\gtrsim 
N(\delta^{-1} I \cap \delta^{-1}J_k)^{1/d}.
$ 
But, by Lemma \ref{idealquotient}, 
\begin{align}\label{I Jk norm calc}
N(\delta^{-1} I \cap \delta^{-1}J_k)^{1/d}
\gtrsim  
N(I \cap J_k)^{1/d}
=(N(I) N(J_k))^{1/d}
\gtrsim 
M_k^{2+\frac{\rho}{d}}, 
\end{align}
where the equality holds because $I$ and $J_k$ are relatively prime. 
Thus the elements of 
$\delta^{-1} I \cap (T^{-1} s + \delta^{-1} J)$ 
%are separated by the same amount.
are separated by 
$\gtrsim 
M_k^{2+\frac{\rho}{d}}.$ 
Since $T$ is a bounded linear isomorphism, 
the elements of $\mathcal{T}(I)$ 
are also separated by $\gtrsim M_k^{2+\frac{\rho}{d}}$. 
It follows that 
\begin{align*}
\sum_{t \in \mathcal{T}(I)} (1+\eta_k|t|)^{-50d}
&\lesssim
\sum_{t \in \ZZ^d} (1+\eta_k M_k^{2+\frac{\rho}{d}}|t|)^{-50d} 
\\
&\lesssim 
\left\{\begin{array}{cc}
    (M_k^{1+\tau} / M_k^{2+\frac{\rho}{d}})^d  & \text{if } \rho \leq d(\tau-1),  \\
    1 & \text{if } \rho > d(\tau-1). 
\end{array}\right.
\end{align*}
In any case, the sum is $\lesssim M_k^{d(\tau-1) -\rho}+1$. 
Plugging this estimate and \eqref{QMsize} into \eqref{h2hatsum2} 
completes the proof. 
\end{proof}

Now we are ready to bound the sum involving $\widehat{h_2}$ in \eqref{phi mu decomp}. 


\begin{lemma}
Suppose $\rho \neq 0$. If $T^{-1} s \in \delta^{-1} J$ and $|s| \leq M_k^{1+\tau}$, then 
\begin{equation}\label{h_2_ub}
\sum_{t \in \ZZ^d} |\widehat{h_2}(s-t) \widehat{\mu_{l,k}}(t)| \lesssim \log^{1/10}(M_k) M_k^{d\tau} (1+ M_k^{\rho -d(\tau-1)}).  
\end{equation}
\end{lemma}




\begin{proof}
By Lemma \ref{h_2_est}, the hypotheses of Lemma \ref{H mu convolution lemma} are satisfied with 
$C_k = M_k^{d \tau} (1 + M_k^{\rho-d(\tau-1)})$ and $H_k = \widehat{h_2}$.  
Consequently, using the assumption $|s| \leq M_k^{1+\tau}$, we get 
\begin{align*}
\sum_{t \in \ZZ^d }|\widehat{h_2} (s-t) \widehat{\mu_{l,k}} (t) | 
\lesssim 
C_k M_k^{-80d(1+\tau)} 
+ 
\sum_{t \in \ZZ^d }
|\widehat{h_2} (s-t) 
\widehat{\mu_{k-1}} (t)|.  
\end{align*}
To bound the sum on the right, we consider separately the ranges $|t| < \log^{1/10d}(M_k)$ and $|t| \geq  \log^{1/10d}(M_k)$.  
In both cases, we shall use the bound  $|\widehat{h}_2(u)| \lesssim  M_k^{d\tau} (1+ M_k^{\rho -d(\tau-1)})$ furnished by Lemma \ref{h_2_est}. 
By \eqref{mu-1}, we have
\begin{align*}
\sum_{|t| < \log^{1/10d}(M_k)} |\widehat{h}_2(s-t) \widehat{\mu}_{k-1}(t)| 
\lesssim  \log^{1/10}(M_k) M_k^{d\tau} (1+ M_k^{\rho -d(\tau-1)}). 
\end{align*}
We assume $M_k$ is large enough that $\log^{1/10d}(M_k) \geq 2M_{k-1}^{1+\tau}$. 
Thus we can apply \eqref{mu-4'} (with $k-1$ in place of $k$) to obtain 
\begin{align*}
\sum_{|t| \geq \log^{1/10d}(M_k)} |\widehat{h}_2(s-t) \widehat{\mu}_{k-1}(t)| 
\lesssim M_k^{d\tau} (1+ M_k^{\rho -d(\tau-1)})\log^{-100}(M_k).
\end{align*}
\end{proof}



We now complete the proof of Lemma \ref{lem_lb_fkmu}. If $\rho \neq 0$, \eqref{f_kmu_l_lb} follows by plugging the estimates \eqref{h_1_lb}, \eqref{h_1_ub}, \eqref{h_2_ub} into 
\ref{phi mu decomp}. 
If $\rho =0$, it suffices to consider $h_1$ only. Indeed, 
\begin{align*}
h_2(x)
= 
c_k (\Phi_{J_k, \eta_k}(x))^2
+ 
\sum_{\substack{I \in Q(M) \\ I \neq J_k}} c_k \Phi_{J_k, \eta_k}(x) \Phi_{I, \eta_k}(x) 
\end{align*}
and the terms of the sum vanish  for $x \in \supp(\phi_0)$ by Lemma \ref{separation lemma}. Hence, if $\rho = 0$, we have
\[h_2(x) = c_k (\Phi_{J_k, \eta_k}(x))^2 = \frac{1}{P_k}  h_1(x) =  h_1(x). \]
So \eqref{f_kmu_l_lb} follows from \eqref{h_1_lb} and \eqref{h_1_ub}.






\subsection{Norm  Bounds}\label{norm bounds}


\begin{lemma}\label{PhiLqnorm}
For each $1 \leq q < \infty$, 
$$
\norm{\Phi_{J_k, \eta_k}}_{L^q(\mu)}^q
\lesssim 
\log^2(M_k) M_k^{qd(1 + \tau) - d + \rho^+}.
$$
\end{lemma}
\begin{proof}
Note $\Phi_{J_k, \eta_k}(x)$ is supported on balls of radius $\eta_k = M_k^{-(1+\tau)}$ which 
(by Lemma \ref{separation facts}) are separated by $\gtrsim M_k^{-1-\frac{\rho}{d}}$. 
As the support of $\mu$ is compact, it intersects $\lesssim M_{k}^{d+\rho}$ of these balls. 
So, by \eqref{mu(B)-3} of Lemma \ref{k-ball regularity}, 
$$
\mu\rbr{\supp\rbr{\Phi_{J_k, \eta_k}}}
\lesssim 
\log^2(M_k)M_{k}^{-d+\rho^+}. 
$$
Since $\Phi_{J_k, \eta_k}$ is bounded above by $\eta_k^{-d} = M_k^{d(1 + \tau)}$, 
the desired result is immediate. 
\end{proof}

\begin{lemma}\label{Phierrorlem}
For all $s \in \ZZ^d$ and all $\xi \in \RR
^d$, 
\begin{equation}\label{Phierroreq}
\left| \widehat{\Phi_{J_k, \eta_k} \mu}(s + \xi) - \widehat{\Phi_{J_k, \eta_k} \mu}(s) \right| \lesssim |\xi| \log^2(M_k) 
M_k^{d \tau+\rho^+}. 
\end{equation}
\end{lemma}
\begin{proof}
Since $\mu$ is supported on 
a compact set, 
we obtain that
\begin{align*}
\left| \widehat{\Phi_{J_k, \eta_k} \mu}(s + \xi) - \widehat{\Phi_{J_k, \eta_k} \mu}(s) \right| & = \left| \int \Phi_{J_k, \eta_k}(x) \left(e^{-2 \pi i (\xi + s) \cdot x} - e^{-2 \pi i s \cdot x} \right) \, d \mu(x) \right|  \\
& \leq \int \Phi_{J_k, \eta_k}(x) 2 \pi |x| |\xi| \, d \mu(x) \\
& \lesssim |\xi| \int \Phi_{J_k, \eta_k}(x) \, d \mu(x). \\
\end{align*}
We apply Lemma \ref{PhiLqnorm} with $q = 1$ to conclude \eqref{Phierroreq}.
\end{proof}

\begin{lemma}\label{count lemma}
Let $S(J_k)$ be the set of $s \in \ZZ^d$ such that $T^{-1} s \in \delta^{-1} J_k$ and $|s| \leq cM_k^{1+\tau}$. The number of elements in $S(J_k)$ is 
$
\gtrsim M_k^{d\tau - \rho}. 
$
\end{lemma}
\begin{proof}
Consider $\cbr{y+J_k : y \in \ZZ^d, \, |y| \leq C_d N(J_k)^{1/d}} \subseteq  \OK/J_k$. 
Here, the constant $C_d$ is chosen so that the number of points $y \in \ZZ^d$ such that $|y| \leq C_d N(J_k)^{1/d}$ is strictly greater than $N(J_k) = |\OK/J_k|$. 
By the pigeonhole principle, there exist distinct $y_1,y_2 \in \ZZ^d$ such that $|y_1|,|y_2| \leq N(J_k)$ and $y_1+J_k = y_2+J_k$. 
Then $y_0 := y_1 - y_2$ is a non-zero element of $J_k$ such that $|y_0| \leq 2C_dN(J_k)^{1/d} \lesssim M_k^{1+\frac{\rho}{d}}$. 
%
Recall that $\ZZ^d \subseteq \delta^{-1}$ and $T(\delta^{-1}) = \ZZ^d$. 
Thus $zy_0 \in \delta^{-1} J_k$ and $T(zy_0) \in \ZZ^d$ for all $z \in \ZZ^d$. 
Moreover, for all $z \in \ZZ^d$, 
$|T(zy_0)|
\lesssim 
|zy_0|
\lesssim 
|z||y_0|
\lesssim |z|M_k^{1+\frac{\rho}{d}}$. 
So there exists a constant $C>0$ such that $|T(zy_0)| \leq c M_k^{1+\tau}$ whenever $z \in \ZZ^d$ and $|z| \leq CM_k^{\tau-\frac{\rho}{d}}$. 
Therefore $T(zy_0) \in S(J_k)$ for all  $z \in \ZZ^d$ with $|z| \leq  CM_k^{\tau - \frac{\rho}{d}}$. 
The number of $z \in \ZZ^d$ with $|z| \leq C M_k^{\tau - \frac{\rho}{d}}$ is  $\gtrsim M_k^{d\tau - \rho}$. 
Since $T$ is invertible and $y_0$ is non-zero, $T(z_1 y_0) \neq T(z_2 y_0)$ whenever $z_1,z_2$ are distinct elements of $\ZZ^d$. 
Therefore the number of elements of $S(J_k)$ is also $\gtrsim M_k^{d\tau - \rho}$. 
\end{proof}


\begin{lemma}\label{PhihatLpnorm v2}
Let $1 \leq p < \infty$. Then 
\[\norm{\widehat{\Phi_{J_k, \eta_k} \mu} }_{L^p}^p \gtrsim \log^{{p - 2d}}(M_k)  
M_k^{(d \tau + \rho^+) p}
M_k^{ (d \tau - \rho)}
\]
\end{lemma}
\begin{proof}
Let $S(J_k)$ be as in Lemma \ref{count lemma}. 
By Lemma \ref{lem_lb_fkmu} and Lemma \ref{Phierrorlem}, if $s \in S(J_k)$ and $|\xi| \leq {1}/{\log^2(M_k)}$, 
then 
\[
\left|\widehat{\Phi_{J_k, \eta_k} \mu}(s + \xi)\right| \gtrsim 
\log(M_k) 
M_k^{d \tau + \rho^+}.
\]
Combining this with Lemma \ref{count lemma} implies the desired result. 
\end{proof}







\begin{proof}[Proof of Proposition \ref{res fail prop}]
By Lemma \ref{PhiLqnorm} and Lemma \ref{PhihatLpnorm v2}, for each odd $k$, 
\begin{align*}
\frac{\norm{\widehat{\Phi_{J_k, \eta_k} \mu}}_{L^p}}{\norm{\Phi_{J_k, \eta_k}}_{L^q(\mu)}} 
\gtrsim  
(\log(M_k))^{\frac{p - 2d}{p} - \frac{2}{q}}
M_k^{\frac{1}{p}(d\tau - \rho)  + \frac{1}{q}(q-1)(\rho^+ - d)}
\end{align*}
If $p < p(\tau, \rho, q, d) := \frac{q(d \tau - \rho)}{(q-1) (d-\rho^+)}$, the exponent of $M_k$ is positive. 
\end{proof}




\section{General Convolution Stability Lemma}\label{convstab section}



\begin{lemma}\label{convstab lemma}
Let $G,H:\ZZ^d \to \CC$ be functions such that $|G| \leq G(0)=1$ and $H$ is bounded.
Let $a > 0$. Let $N$ be a positive integer such that 
$N>2d+2a$. Let $1 \leq A \leq B$. 
Let $D$ be a number satisfying 
\begin{equation}\label{Dcond}
B^{-(N - d)} \leq A^{-(N-d)} \leq D
\end{equation} 
Let $L: (1,\infty) \to (0,\infty)$ be a function such that, for all sufficiently large $x$, 
\begin{align}\label{Mcond6}
\text{$L(x)$ is increasing 
}
\end{align}
%
\begin{align}\label{Mcond2} 
\text{$L(x)/x^a$ is decreasing } 
\end{align}


Suppose the following estimates hold for all $s \in \ZZ^d$: 
\begin{align}
\label{convcond1} 
G(s) = 0 \quad \text{if } 1 \leq |s| \leq A; 
\end{align}
\begin{align}
\label{convcond2}
|G(s)| \lesssim D L(|s|) \quad \text{if } |s| \geq A;  
\end{align}
\begin{align}
\label{convcond3}
|G(s)| \lesssim D (B / |s|)^N L(|s|)  \quad \text{if } |s| \geq B;  
\end{align}
\begin{align}
\label{convcond4}
|H(s)| \lesssim |s|^{-(N-d)} L(|s|)  \quad \text{if } |s| \geq 1. 
\end{align}
%
If $A$ and $B$ 
are sufficiently large 
depending 
on 
the other parameters  
and on the implicit constants in \eqref{convcond2}, \eqref{convcond3}, and \eqref{convcond4}, then the following estimates hold for all $s \in \ZZ^d$:
%
\begin{align}\label{convlemmares1}
|G \ast H(s) - H(s)| 
\lesssim 
A^{-(N-2d)} L(A/2)
%\leq
%A^{-(N-2d-a)} L(A/2)
\quad \text{if } |s| \leq A/2; 
\end{align}
\begin{align}\label{convlemmares2}
|G \ast H(s)| %\leq  
\lesssim 
D 
%L(A) 
L(2|s|)
\quad \text{if } |s| \geq A/2 ; 
\end{align}
\begin{align}\label{convlemmares3}
|G \ast H(s)| 
%\leq 
\lesssim 
D (B/|s|)^{N-d} %L(A) 
L(|s|/2)
\quad \text{if } |s| \geq 2 B. 
\end{align}
\end{lemma}




\begin{proof}[Proof of \eqref{convlemmares1}] 
Assume $|s| \leq \dfrac{1}{2} A$. 
By \eqref{convcond1} and because  
$|G| \leq G(0)=1$, 
we have 
$$
|G \ast H(s) - H(s)| 
= \abs{\sum_{|t| \geq A} G(t) H(s-t)}
\leq \sum_{|t| \geq A} |H(s-t)|. 
$$
The assumption $|s| \leq \dfrac{1}{2} A$ and the condition $|t| \geq A$ on the sum imply that 
$|s-t| \geq A/2$. 
Thus 
$$
|G \ast H(s) - H(s)| \leq \sum_{ |s-t| \geq A/2 } |H(s-t)|. 
$$
Now \eqref{convcond4} implies 
\begin{align*}
|G \ast H(s) - H(s)| 
&\lesssim 
\sum_{ |s-t| \geq A/2  }|s-t|^{-(N-a-d)} \dfrac{L(|s-t|)}{|s-t|^a}.  
\end{align*} 
By \eqref{Mcond2}, assuming $A$ is sufficiently large, 
$L(x)/x^a$ is decreasing for all $x \geq \dfrac{1}{2} A$. 
Moreover, $N - a - d > d$ by assumption. 
Therefore 
$$
|G \ast H(s) - H(s)| 
\lesssim 
A^{-(N-2d)} L(A/2). 
$$
%
\end{proof}



\begin{proof}[Proof of \eqref{convlemmares2}] 
Assume $|s| \geq A/2$. Assume $A$ is large enough that 
$L(x)$ is increasing 
and $L(x)/x^a$ is decreasing for all $x \geq A/2$. 
By \eqref{convcond1} and because $G(0)=1$, we have 
$$
G \ast H(s) = H(s) + S_1 + S_2 %H(s) + \sum_{|s-t| \geq A} G(s-t) H(t),
$$
where 
\begin{align*}
S_1 &= \sum_{\substack{|s-t| \geq A \\ |t| \leq |s|}} G(s-t) H(t), \\
S_2 &= \sum_{\substack{|s-t| \geq A \\ |t| \geq |s|}} G(s-t) H(t).
\end{align*}
%
By  \eqref{convcond4}, 
$$
|H(s)| \lesssim  |s|^{-(N-d)} L(|s|) \lesssim 
A^{-(N-d)} L(2|s|). 
$$
%
By the assumption that 
$A^{-(N-d)} \leq D$, we have 
$$|H(s)| \lesssim  DL(2|s|).$$
%
By \eqref{convcond2}, 
$$
|S_1| \lesssim 
D L(2|s|) \sum_{\substack{|s-t| \geq A \\ |t| \leq |s|}} |H(t)|
\leq 
D L(2|s|) \sum_{t \in \ZZ^d} |H(t)|.  
$$
Because $H$ is bounded and because of \eqref{Mcond2} and \eqref{convcond4}, 
the last sum is $\lesssim 1$. 
Thus 
$$
|S_1| \lesssim D L(2|s|). 
$$

On the other hand, by \eqref{convcond2} and \eqref{convcond4}, 
$$
|S_2| \lesssim 
D \sum_{|t| \geq |s|} |t|^{-(N-d-2a)} \dfrac{L(2|t|)}{(2|t|)^a} \dfrac{L(|t|)}{|t|^a}. 
$$
By \eqref{Mcond2} and because $N-d-2a>d$ by assumption, 
$$
|S_2| \lesssim 
D |s|^{-(N-2d)} L(2|s|) L(|s|)  
\lesssim 
D |A|^{-(N-2d)} L(2|s|) L(A). 
$$
%
Using \eqref{Mcond2} and $N-d-2a>d$ again, if $A$ is large enough, we have 
$$
|S_2| \leq DL(2|s|).
$$
\end{proof}



 

\begin{proof}[Proof of \eqref{convlemmares3}] 
Assume $|s| \geq 2B$. Assume $A$ is large enough that $L(x)$ is increasing and $L(x)/x^a$ is decreasing for all $x \geq A$. 
By \eqref{convcond1} and because $G(0)=1$, we have 
$$
G \ast H(s) = H(s) + S_1 + S_2 + S_3,
$$
where 
\begin{align*}
S_1 &= \sum_{\substack{A \leq |t| < B \\ |s - t| \geq |s|/2}} G(t) H(s-t),  
\\
S_2 &= \sum_{\substack{|t| \geq |s|/2 \\ |s - t| \leq |s|/2}} G(t) H(s-t), \\
S_3 &= \sum_{\substack{|t| \geq B \\ |s - t| \geq |s|/2}} G(t) H(s-t). 
\end{align*}
By \eqref{convcond4}, 
$$
|H(s)| \lesssim |s|^{-(N-d)} L(|s|) \leq (|s|/2)^{-(N-d)} L(|s|/2). 
$$
%
Because of the assumption $B^{-(N-d)} \leq D$, we have 
$$
|H(s)| \lesssim D (B/|s|)^{N-d} L(|s|/2). 
$$
%
By \eqref{convcond2} and \eqref{convcond4}, 
\begin{align*}
|S_1| 
&\lesssim |s|^{-(N-d)}L(|s|/2) \sum_{\substack{A \leq |t| < B \\ |s - t| \geq |s|/2}} |G(t)|
\\
&\lesssim |s|^{-(N-d)}L(|s|/2) 
\sum_{\substack{A \leq |t| < B \\ |s - t| \geq |s|/2}} DL(|t|) 
\\
&\lesssim 
D |s|^{-(N-d)} L(|s|/2) L(B) B^d
\end{align*}
Since $L(x)/x^a$ is decreasing for $x \geq A$ and since $N-d>d+a$, we have 
$L(B) B^d \leq L(A) A^{-a} B^{N-d}$. 
%%\begin{comment}
Moreover, by 
\eqref{Mcond2}, $L(A) A^{-a} \lesssim 1$.
%%\end{comment}
Thus 
$$
|S_1| \lesssim  D (B/|s|)^{N-d} L(|s|/2).
$$
By \eqref{convcond3},  
\begin{align*}
|S_2| &\lesssim D(B/|s|)^N L(|s|/2)
\sum_{\substack{|t| \geq |s|/2 \\ |s - t| \leq |s|/2}}  |H(s-t)| 
\\
&\leq 
D(B/|s|)^N L(|s|/2) \sum_{t \in \ZZ^d}  |H(s-t)|.
\end{align*}
Because $H$ is bounded, because $L(x)/x^a$ is eventually decreasing, and because of \eqref{convcond4}, the last sum is bounded by a constant independent of $s$. 
Thus 
$$
|S_2| \lesssim D (B/|s|)^{N-d} L(|s|/2). 
$$
By \eqref{convcond3} and \eqref{convcond4}, 
\begin{align*} 
|S_3| 
&\lesssim 
|s|^{-(N - d)} L(|s|/2) \sum_{\substack{|t| \geq B \\ |s - t| \geq |s|/2}} |G(t)| \\
&\lesssim 
D B^N |s|^{-(N - d)} L(|s|/2) \sum_{\substack{|t| \geq B \\ |s - t| \geq |s|/2}}  |t|^{-(N-a)} \dfrac{L(|t|)}{|t|^a}. 
\end{align*}
Since $N-a>d$ and since $L(x)/x^a$ is decreasing for $x \geq B$, the last sum is $\lesssim B^{-(N-d)}L(B)$. Therefore 
$$
|S_3| 
\lesssim 
D |s|^{-(N - d)} L(|s|/2) L(B) B^d. 
$$
Since $L(x)/x^a$ is decreasing for $x \geq A$ and since $N-d>d+a$, we have 
$L(B) B^d \leq L(A)A^{-a}B^{N-d}$. 
Moreover, by \eqref{Mcond2}, $L(A) A^{-a} \lesssim 1$. 
Thus 
$$|S_3| \lesssim D (B/|s|)^{N-d} L(|s|/2).$$
\end{proof}









\section{Hausdorff Dimension Upper Bound}\label{Hausdorff Dimension Upper Bound section}


\begin{prop}\label{hausdorff upper bound}
$\dim_H(E(K,B,\tau)) \leq 2d/(1+\tau)$.  
\end{prop}
%
\begin{proof}
Let $\mathcal{I}(\OK)$ denote the set of all ideals of $\OK$. 
Note $E(K,B,\tau)$ is invariant under translation by elements of $\ZZ^d$ because for $\ZZ^d \isom \OK \subseteq I^{-1}$ for every $I \in \mathcal{I}(\OK)$. 
Thus it suffices to show that $E(K,B,\tau) \cap [-1/2,1/2)^d$ has Hausdorff dimension at most $2d/(1+\tau)$. 
Let $B(r,M)$ denote the closed ball in $\RR^d$ with center $r$ and radius $M$. 
Note 
\begin{align*}
E(K,B,\tau) 
 = \bigcap_{n=1}^{\infty} \bigcup_{\substack{ I \in \mathcal{I}(\OK) \\ N(I) > n}}  \bigcup_{r \in I^{-1}} 
 B( r,N(I)^{-(1+\tau)/d} ). 
\end{align*}
So, for every $n \geq 1$,  
\begin{align*}
E(K,B,\tau) \cap [-1/2,1/2)^d 
\subseteq \bigcup_{\substack{ I \in \mathcal{I}(\OK) \\ N(I) > n}} 
\bigcup_{\substack{r \in I^{-1} \\ |r| < 2}} 
 B( r,N(I)^{-(1+\tau)/d} ). 
\end{align*}
Since $\OK \isom \ZZ^d$, for every $a \in \RR^d$, 
the set $I^{-1} \cap (a+[0,1)^d)$ is a complete set of representatives for $I^{-1}/\OK$. Thus the set of $r \in I^{-1}$ with $|r|<2$ is covered by at most $4^d$ complete sets of representatives of $I^{-1}/\OK$. By Lemma \ref{idealquotient}, 
$|I^{-1} / \OK| 
= {N(\OK)}/{N(I^{-1})} 
= N(I)$.  
Therefore the set of $r \in I^{-1}$ such that $|r| < 2$ has size $\lesssim N(I)$. Thus, for every $s \geq 0$, 
\begin{align*}
\sum_{\substack{ I \in \mathcal{I}(\OK) \\ N(I) > n}} 
\sum_{\substack{r \in I^{-1}\\ |r| < 2}} 
	\rbr{ \text{diam}( B( r,N(I)^{-(1+\tau)/d} ) ) }^s
\lesssim   
\sum_{\substack{ I \in \mathcal{I}(\OK) \\ N(I) > n}} N(I)^{1-s(1+\tau)/d}. 
\end{align*}
If $s > 2d/(1+\tau)$, the last sum goes to zero as $n \to \infty$. (This is because the sum $\zeta_{K}(z) := \sum_{ I \in \mathcal{I}(\OK) } N(I)^{-z}$, known as the Dedekind zeta function, converges whenever $\Re(z)>1$; see \cite[Ch.10]{Jarvis14} for a proof.)
Therefore the Hausdorff dimension of $E(K,B,\tau) \cap [-1/2,1/2]^d$ is at most $2d/(1+\tau)$. 
\end{proof}

\section{Appendix: modifications for \texorpdfstring{$E_{\prin}(K,B,\tau)$}{Eprin(K,B,tau)}}\label{sec_prin}
We will describe the modifications necessary to prove that Theorem \ref{mainthm3} holds for the set $E_{\prin}(K,B,\tau)$. The key 
%to adapting the argument to apply to $E_{\prin}$ 
is to modify the argument to apply to \textit{principal} ideals. This does not change the argument much because of the following standard fact from algebraic number theory.
\begin{thm}[Finiteness of the ideal class group]\label{idealclassthm}
Let $K$ be an algebraic number field. Then there exists a finite collection $\mathcal{H}$ of integral ideals such that if $I$ is an ideal of $\OK$, then there exists $H_I \in \mathcal{H}$ such that the ideal $I H_I$ is a principal ideal.
\end{thm}

Using Theorem \ref{idealclassthm}, we define the following analog of $Q(M)$ consisting of principal ideals: 
\[Q_{\prin}(M) = \left\{I H_I : I \in Q(M/R) \right\}\]
where 
$$
R = \max\cbr{ N((H_{1}^{-1} + H_{2}^{-1})^{-1}) : H_1, H_2 \in \mathcal{H} }. 
$$
%
%
Then $Q_{\prin}(M)$ is a collection of $\sim {M^d}/{\log M}$ principal ideals of norm $\sim M^d$ that are ``almost" prime ideals. 

Everywhere in the original argument, we replace $Q(M)$ by $Q_{\prin}(M)$.   
%We also replace $J_k$ with the  principal ideal $J_k H_{J_k}$. 
In particular, this means $J_k \in Q_{\prin}(M_k)$. 

Lemma \ref{primedivisorbound} is replaced by the following lemma (which is actually a corollary of Lemma \ref{primedivisorbound}). 

\begin{lemma}\label{principaldivisorbound}
Let $J$ be a nonzero integral ideal of $\OK$. Then 
\begin{equation}\label{principaldivisorineq}
|\{I \in Q_{\prin}(M) : J \subseteq I \}| \leq \frac{\log (2N(J))}{d \log (M/R)}.
\end{equation}
\end{lemma}
\begin{proof}
This follows from Lemma \ref{primedivisorbound} upon observing that 
\begin{align*}
|\cbr{I \in Q_{\text{prin}}(M) : J \subseteq I    }|
&=
|\cbr{I H_I : I \in Q(M/R), J \subseteq I H_I}| 
\\
&\leq 
|\cbr{I H_I : I \in Q(M/R), J \subseteq I}| 
\\
&\leq 
|\cbr{I : I \in Q(M/R), J \subseteq I}|.
\end{align*}
\end{proof}

The proof of Lemma \ref{separation facts}(c) must be changed because distinct ideals in $Q_{\prin}(M_k)$ may not be relatively prime. 
To deal with this, we modify the definition of $\phi_0$ so that its support is contained in $([-3/(8R^{1/d}), -1/(8R^{1/d})] \cap [1/(8R^{1/d}), 3/(8R^{1/d})])^d$. Here is a restatement of Lemma \ref{separation facts}(c) and its new proof. 
\begin{lemma}
Let $k \in \mathbb{N}$. Suppose $I_1, I_2 \in Q_{\prin}(M_k)$, $r_1 \in I_1^{-1}, I_{2}^{-1}$. 
Then $I_1 = I_1' H_{I_1'}$ and $I_2 = I_2' H_{I_2'}$ for some $I_1',I_2' \in Q(M_k/R)$.  
If $r_1 = r_2$ and $I_1 \neq I_2$, then 
$r_1 = r_2 \in H_{I_1'}^{-1} + H_{I_2'}^{-1}$
 and (consequently) 
$\dist(r_1,\supp(\phi_0)) = \dist(r_2,\supp(\phi_0)) > 3\eta_k$. 
\end{lemma}
\begin{proof}
Since $I_1 \neq I_2$, we must have $I_1' \neq I_2'$.  
Thus $I_1'$ and $I_2'$ are distinct prime ideals and (hence) are relatively prime, i.e., $1 \in I_1'+I_2'$. 
Set $r = r_1 = r_2$. 
Since $r \in I_1^{-1}$ and $r \in I_2^{-1}$, we have $r I_1' \subseteq H_{I_1'}^{-1}$ and $r I_2' \subseteq H_{I_2'}^{-1}$. 
Therefore $r = r \cdot 1 \in r(I_1' + I_2') = r I_1' + r I_2' \subseteq H_{I_1'}^{-1} + H_{I_2'}^{-1}$.
%%
Thus, if $r \neq 0$, then 
$$
|r|^d \gtrsim N(r) \geq N(H_{I_1'}^{-1} + H_{I_2'}^{-1}) \geq R^{-1}. 
$$
%%

%%
Since $\supp(\phi_0)$ is contained in $([-3/(8R^{1/d}), -1/(8R^{1/d})] \cap [1/(8R^{1/d}), 3/(8R^{1/d})])^d$, we get the desired result. 
\end{proof} 
%
%
Note that (with the modifications to $\phi_0$)  
all the estimates on $F_k$, $\mu_k$, and the other functions and measures still hold with suitable modifications to the implied constants (which now depend on $R$). 
%

Finally, the calculation  \eqref{I Jk norm calc} in the proof of Lemma \ref{h_2_est} must be modified. 
Now $I = I' H_{I'}$ and $J_k = J_k' H_{J_k'}$ 
for some $I' \in Q(M_k/R)$ and 
$J_k' \in Q(M_k^{1+\frac{\rho}{d}}/R)$. 
So \eqref{I Jk norm calc} becomes 
\begin{align*}
N(\delta^{-1} I \cap \delta^{-1} J_k)^{1/d}
&=
N(\delta^{-1} I' H_{I'} \cap \delta^{-1} J_k' H_{J_k'})^{1/d} 
\gtrsim 
N(I' H_{I'} \cap J_k' H_{J_k'})^{1/d}
\\
&\geq  
N(I' \cap J_k')^{1/d} 
= 
(N(I')N(J_k'))^{1/d}
\gtrsim M_k^{2+\frac{\rho}{d}}.
\end{align*}
No other changes are needed. 



%\section*{Acknowledgements}







\bibliographystyle{myplain}
\bibliography{Approximation_In_Number_Fields}
\end{document}



